\documentclass[11pt]{article}

\usepackage[margin=1.1in]{geometry}
\usepackage{amsmath}
\usepackage{amssymb}
\usepackage{amsthm}
\usepackage{booktabs}
\usepackage{array}
\usepackage{longtable}
\usepackage[T1]{fontenc}
\usepackage[htt]{hyphenat}
\usepackage[hidelinks]{hyperref}

\theoremstyle{plain}
\newtheorem{theorem}{Theorem}[section]
\newtheorem{lemma}[theorem]{Lemma}
\newtheorem{proposition}[theorem]{Proposition}
\newtheorem{corollary}[theorem]{Corollary}
\theoremstyle{definition}
\newtheorem{definition}[theorem]{Definition}
\newtheorem{problem}[theorem]{Problem}
\newtheorem{question}[theorem]{Question}
\theoremstyle{remark}
\newtheorem{remark}[theorem]{Remark}

\numberwithin{equation}{section}

\newcommand{\R}{\mathbb{R}}
\newcommand{\rad}{\operatorname{rad}}
\newcommand{\ecc}{\operatorname{ecc}}
\newcommand{\diam}{\operatorname{diam}}
\newcommand{\pth}{\operatorname{path}}
\newcommand{\comp}{\operatorname{comp}}
\newcommand{\ball}{\operatorname{ball}}
\newcommand{\gir}{\operatorname{girth}}

% ---- verification-status markers -------------------------------------------
\emergencystretch=4em
\hbadness=10000
\newcommand{\statusbox}[1]{{\small\sloppy\textsf{[#1]}}}
\newcommand{\PRV}[1]{\statusbox{proved here; assertions machine-checked by \texttt{#1}}}
\newcommand{\PRVn}{\statusbox{proved here}}
\newcommand{\REV}[1]{\statusbox{proved here; machine-checked by \texttt{#1};
\textbf{two-family review}}}
\newcommand{\AD}[1]{\statusbox{machine-assisted proof, re-derived here;
machine-checked by \texttt{#1}; \textbf{one adversarial review pass}}}
\newcommand{\MV}[1]{\statusbox{machine-verified: \texttt{#1}}}
% ----------------------------------------------------------------------------

\title{Partial Results on the WOWII-133 Graffiti Conjecture\\
{\large The $C_4$-free branch below $\lfloor \ell \rfloor = 4$}}

\author{Haoyu Chen\thanks{%
\textbf{Disclosure of the production pipeline.}
This paper reports the certified by-products of an autonomous, still-running attack on
Conjecture~133 of the \emph{Written on the Wall II} list. \emph{The conjecture is not
resolved here, and this paper does not claim to resolve it.} Every mathematical statement
below was produced by automated systems: an orchestrating agent (Claude Fable~5) that set
the attack targets, adjudicated the returning reports line by line, executed the promotion
protocol of \S\ref{sec:conventions} and drafted this document; a free exploration channel
(Qwen3.8-Max, web interface) that was given self-contained task briefs and whose output was
never adopted without independent re-derivation; adversarial reviewers drawn from a model
family disjoint from the one that produced the argument under review (Claude Opus-family
instances and Qwen3.8-Max, each run blind or diff-scoped, in isolation, forbidden to consult
the project's other files); and machine checkers (purpose-written C and Python enumerators
and independent from-scratch re-verifiers). No human supplied a definition, a lemma, a proof
idea, or a correction. Provenance is recorded per result in \S\ref{sec:provenance}: several
of the arguments below originated in the free channel and were re-proved before being
admitted, and one of the paper's results --- Theorem~\ref{thm:g45}, a refutation --- was
found by a reviewer whose assignment was to look for defects. Several of the campaign's own
working claims were refuted by its own later rounds; those self-refutations are recorded in
\S\ref{sec:limits} rather than quietly deleted. Responsibility for what follows rests with
the reader who re-runs the scripts of Appendix~\ref{app:artifacts}.}}

\date{22 August 2026}

\begin{document}

\maketitle

\begin{abstract}
Conjecture~133 of the \emph{Written on the Wall II} series, produced by the Graffiti.pc
conjecture-making program, asserts that a connected graph $G$ satisfies
$\rad(G) + \lfloor \ell(G)\rfloor^{c} \le \pth(G)$, where $\pth(G)$ is the order of a
longest induced path, $\ell(G)$ is the average over $v$ of the independence number of the
open neighbourhood $N(v)$, and $c = 1$ if $G$ is $C_4$-free and $c = 0$ otherwise. We report
partial progress. The $C_4$-containing branch and the whole range $\lfloor \ell\rfloor \le 2$
are settled unconditionally. For the first open band $\lfloor \ell\rfloor = 3$ we prove that
a counterexample of minimum order must have radius at least $4$: the radius-$2$ layer is
closed by a fibre-counting theorem, and the radius-$3$ layer is closed by a rigidity
analysis of the induced $6$-cycle that such a counterexample is forced to carry, whose engine
is an unconditional rigidity lemma --- in a $C_4$-free graph with no induced $P_6$, a vertex
outside an induced $C_6$ has $0$ or $2$ neighbours on the cycle, and if $2$ they are
antipodal. At radius $\ge 5$ in the band $\lfloor \ell\rfloor \ge 4$ we close the case
$\mu \ge 2$ and reduce what remains to a single configuration, which we are careful to note
is not exhibited by any graph we know; a $54$-vertex witness refutes the reformulation of that
residual which we had been treating as equivalent to it. We also prove the $D = 2$ case of the
natural strengthening
$\pth \ge \diam + 4$ for $C_4$-free graphs of diameter $D$ with $\ell > 3$, and, in the band
$\lfloor \ell\rfloor \ge 4$, that a $C_4$-free graph all of whose peripheral geodesics have
low-independence side-neighbours satisfies $\ell < 4$. A deliberate part of the report is
negative: we record an infinite family showing that no two-endpoint local extension lemma
can reach $\rad + k$ for $k \ge 4$, a $150$-vertex witness killing the layered-counting route
at radius $3$, and two counterexamples showing that the deep endpoint-extension lemma
available at geodesic length $d \ge 4$ is false at $d = 3$ both with and without the natural
repair hypothesis. Every statement carries an explicit epistemic marker, and the verification
scripts are inventoried in an appendix.
\end{abstract}

\tableofcontents

\section{Introduction}\label{sec:intro}

\subsection{The conjecture}

Throughout, $G$ is a finite simple graph on $n = |V(G)| \ge 2$ vertices with $m$ edges and
$T$ triangles, $d(v)$ is the degree of $v$, $t(v)$ the number of triangles through $v$,
$\delta$ and $\Delta$ the minimum and maximum degree, $\rad$, $\diam$, $\ecc$ the radius,
diameter and eccentricity, and $N(v)$ the \emph{open} neighbourhood of $v$. We write
\[
  a(v) := \alpha\bigl(G[N(v)]\bigr), \qquad
  \ell(G) := \frac1n \sum_{v \in V} a(v), \qquad
  \mu(G) := \min_{v} a(v),
\]
where $\alpha$ is the independence number, and $\pth(G)$ for the number of \emph{vertices} of
a longest induced (chordless) path of $G$. We call $a(v)$ the \emph{$a$-value} of $v$ and
$\ell$ the \emph{neighbourhood independence average}.

We take the statement of the conjecture from its Lean transcription in the
\texttt{formal-conjectures} corpus, file
\texttt{FormalConjectures/} \texttt{WrittenOnTheWallII/} \texttt{GraphConjecture133.lean},
where it is recorded as conjecture $133$ of the \emph{Written on the Wall II} series produced by the Graffiti.pc
conjecture-making program of E.~DeLaVi\~na. We have not consulted the original \emph{Written
on the Wall II} document; the numbering and the attribution here are those of the Lean file,
and the reader who needs the primary source should check it. The Lean statement reads

\begin{verbatim}
theorem conjecture133 (G : SimpleGraph a) [DecidableRel G.Adj]
    (h : G.Connected) :
    let rad := G.radius.toNat
    let hasC4 := exists a b c d, (pairwise distinct) and
      G.Adj a b and G.Adj b c and G.Adj c d and G.Adj d a
    let cC4 : N := if hasC4 then 0 else 1
    (rad : R) + (floor (l G) : R) ^ cC4 <= (path G : R)
\end{verbatim}

Three points of the transcription matter and were checked against the Lean sources before any
mathematics was attempted. First, \texttt{path} is the induced-path order, not the
``average distance'' homonym that the same corpus uses elsewhere. Second, \texttt{hasC4} asks
for a $4$-cycle as a \emph{subgraph}, chords allowed; equivalently, $G$ fails to be
$C_4$-free exactly when some two distinct vertices have two common neighbours. This
\emph{strong} form of $C_4$-freeness is the one used everywhere below, and the distinction is
load-bearing: several lemmas are false under the induced-only reading. Third, the exponent
$c \in \{0,1\}$ is applied through \texttt{Monoid.npow}, so on the $C_4$-containing branch the
right-hand side is $\rad + 1$ whatever $\ell$ is.

Both sides being integers cast into $\R$, the Lean goal is equivalent to the conjunction of

\begin{itemize}
\item \textbf{(C133-a)} \; $G$ connected with a $4$-cycle: \; $\pth(G) \ge \rad(G) + 1$;
\item \textbf{(C133-b)} \; $G$ connected and $C_4$-free: \;
      $\pth(G) \ge \rad(G) + \lfloor \ell(G)\rfloor$.
\end{itemize}

\subsection{What is proved here, and what is not}

The conjecture is open. Table~\ref{tab:status} is the honest ledger; the numbers in it are
justified in the sections named.

\begin{table}[ht]
\centering
\small
\begin{tabular}{@{}lll@{}}
\toprule
Region & Status & Where \\
\midrule
$G$ contains a $4$-cycle & closed & Thm.~\ref{thm:A} \\
$C_4$-free, $\lfloor\ell\rfloor \le 1$ & closed & Thm.~\ref{thm:B} \\
$C_4$-free, $\lfloor\ell\rfloor = 2$ & closed & Thm.~\ref{thm:D}, Cor.~\ref{cor:D1} \\
$C_4$-free, $\lfloor\ell\rfloor = 3$, $\mu \ge 2$ & closed & Thm.~\ref{thm:gplus},
   Cor.~\ref{cor:gplus1} \\
$C_4$-free, $\lfloor\ell\rfloor = 3$, $\rad = 2$ & closed & Thm.~\ref{thm:g7},
   Cor.~\ref{cor:g8} \\
$C_4$-free, $\lfloor\ell\rfloor = 3$, $\rad = 3$ & closed \emph{for a minimum-order
   counterexample} & Thm.~\ref{thm:r3} \\
$C_4$-free, $\lfloor\ell\rfloor = 3$, $\rad \ge 4$ & \textbf{open}; reduced to
   Problem~\ref{prob:rpd} & \S\ref{sec:peeling} \\
$C_4$-free, $\lfloor\ell\rfloor \ge 4$, $\rad \ge 5$, $\mu \ge 2$ & closed &
   Thm.~\ref{thm:g42} \\
$C_4$-free, $\lfloor\ell\rfloor \ge 4$, otherwise & \textbf{open}; partial results in
   \S\ref{sec:adopted}, \S\ref{sec:pocket2} & Table~\ref{tab:pocket2} \\
\bottomrule
\end{tabular}
\caption{Status of Conjecture~133 after the work reported here.}
\label{tab:status}
\end{table}

The main contributions are these.

\begin{enumerate}
\item \textbf{The closed half, with sharp constants.} Theorem~\ref{thm:A} settles (C133-a) in
four lines. Theorem~\ref{thm:D} --- every connected $C_4$-free graph of radius $\ge 2$ has
$\pth \ge \rad + 2$ --- settles $\lfloor \ell\rfloor \le 2$. Both are tight on the five
extremal graphs $K_2, K_3, C_5, C_6$ and Petersen (\S\ref{sec:falsification}), which is what
forbids any argument that loses a constant.

\item \textbf{An endpoint-extension calculus and the band $\lfloor\ell\rfloor = 3$ with no
$a = 1$ vertex.} Lemmas~\ref{lem:g10}, \ref{lem:g1}, \ref{lem:g9} extend a geodesic by
one or two vertices from data at its ends. Theorem~\ref{thm:gplus} combines them:
$\mu \ge 2$ together with one vertex of $a$-value $\ge 3$ forces $\pth \ge \rad + 3$. This
alone closes every $C_4$-free graph on at most $13$ vertices (\S\ref{sec:falsification}).

\item \textbf{The radius-$2$ layer.} Theorem~\ref{thm:g7}: a connected $C_4$-free graph with
$\rad = 2$ and $\ell \ge 3$ always contains a pair of vertices at distance $2$ with
$a$-values $(\ge 3, \ge 2)$. The proof is a fibre decomposition of the ball of radius $2$
around a high vertex followed by an exact count. Together with Lemma~\ref{lem:g1} this closes
$\rad = 2$.

\item \textbf{The radius-$3$ layer} (\S\ref{sec:r3}, the longest part of the paper). A
minimum-order counterexample of radius $3$ is reduced --- by a peeling induction whose
residual is made explicit --- to a graph $C$ that is $C_4$-free, has $\rad(C) = 2$,
$\diam(C) = 3$, $\ell(C) > 3$, no induced $P_6$, and the extra property (R3b) that every
vertex of $a$-value $\ge 3$ has eccentricity $\le 2$. Two theorems empty that class.
Theorem~\ref{thm:g26} shows no vertex of such a graph has $a$-value $1$ and eccentricity $3$;
Theorem~\ref{thm:g35} then shows the surviving branch is empty, by way of
Lemma~\ref{lem:g32}, an unconditional rigidity fact about induced $6$-cycles in $P_6$-free
$C_4$-free graphs which collapses the configuration to at most nine vertices.

\item \textbf{The integer form of the floor-free strengthening at $D = 2$.}
Theorem~\ref{thm:g12}: $C_4$-free, $\diam = 2$, $\ell > 3$ imply $\pth \ge 6 = \rad + 4$.
This is a non-trivial instance of the real-valued strengthening
$\ell \le \pth - \rad$ (\S\ref{sec:aux}) and is the $D = 2$ case of Problem~\ref{prob:rpd},
to which all of $\rad \ge 4$ has been reduced.

\item \textbf{The band $\lfloor\ell\rfloor \ge 4$} (\S\ref{sec:adopted}). Three results at the
lower reviewed tier: the deep endpoint-extension lemma known at geodesic length $d \ge 5$ is
extended to $d = 4$ unconditionally (Lemma~\ref{lem:g36}) and to $d = 3$ under one extra
hypothesis (Lemma~\ref{lem:g37}); and Theorem~\ref{thm:g38}, a global counting theorem: a
connected $C_4$-free graph in which every geodesic of length $\ge \rad$ has all its usable
side-neighbours of $a$-value $\le 3$ satisfies $\ell < 4$. Contrapositively, every $C_4$-free
graph with $\ell > 4$ supplies, at every radius, a peripheral geodesic with a side-neighbour
of $a$-value $\ge 4$.

\item \textbf{The same band at radius $\ge 5$} (\S\ref{sec:pocket2}), at the higher
two-family reviewed tier. Theorem~\ref{thm:g42} closes the slice $\rad \ge 5$,
$\mu \ge 2$ outright, and Lemma~\ref{lem:g43} with Corollary~\ref{cor:g44} reduces what
remains at $\rad \ge 5$ to a single configuration --- one that no known graph realises, a
limit on the reduction which Remark~\ref{rem:nonexhibited} states rather than files away.
Proposition~\ref{prop:w2} supplies the regime's hard positive control, a $660$-vertex Cayley
graph of $\mathrm{PSL}(2,11)$; Theorem~\ref{thm:g48} refutes, by a $54$-vertex witness, the
reformulation of the residual that the campaign had been treating as equivalent, and kills
with the same graph the assumption that peeling preserves the radius.

\item \textbf{Negative results} (\S\ref{sec:limits}), reported because they close routes.
An infinite family $T_m$ of trees shows that ``both endpoints of $a$-value $\ge k$ implies
$\pth \ge \rad + k$'' is false for every $k \ge 4$. A $150$-vertex graph (Hoffman--Singleton
with two pendant vertices at each vertex) has $\lfloor\ell\rfloor = 3$, radius $3$, and
\emph{no} pair at distance $\ge \rad$ with $a$-values $(\ge 3, \ge 2)$, which kills the
layered-counting route above radius $2$. Two ten-vertex graphs refute the $d = 3$ analogue of
Lemma~\ref{lem:g36}, one without and one with the natural repair hypothesis
(Theorem~\ref{thm:g45}); the conclusion is that the free local route at $d = 3$ is closed and
a global hypothesis must be paid for. Finally we record two defects that the review protocol
found in the project's own earlier text --- a statement stated more widely than its proof
supported, and an arithmetic error in the peeling bookkeeping --- together with their
repairs, since both are instructive.
\end{enumerate}

Nothing here is formalised in Lean, and no claim below should be read as a formal
verification. The scripts of Appendix~\ref{app:artifacts} check assertions over stated finite
ranges and on named explicit graphs; the proofs are ordinary mathematics.

\subsection{Verification conventions and epistemic markers}\label{sec:conventions}

Because this is an automated campaign's output, the status of each statement is recorded
explicitly and is \emph{not} uniform. Four markers are used, attached to every numbered
assertion; definitions and open problems assert nothing and carry none.

\begin{itemize}
\item \PRV{script} --- the proof is reproduced here in full, and the assertions it makes are
checked by the named script over exactly the range stated at the point of use, usually an
exhaustive table of small graphs plus explicit large witnesses. A script never checks more
than the statement it is attached to.

\item \REV{script} --- as above, and in addition the statement passed the campaign's
\emph{two-family} adversarial protocol: an independent reviewer from a model family disjoint
from the one that produced the argument, run blind (given the text and nothing else, with no
access to the project's other files and no internet, permitted to write and run its own
verification code), \emph{plus} a second, diff-scoped pass over the same material by a
different reviewer, both returning no surviving defect at any joint classified as
mathematics. Every such statement in this paper survived at least one reviewer finding a
\emph{real} defect elsewhere in its neighbourhood; the repairs are in the text.

\item \AD{script} --- \emph{machine-assisted}: the argument originated in an automated solve
attempt in the free exploration channel, was re-derived line by line here before admission
(and in two places repaired --- the repairs are marked in the proofs), is machine-checked as
above, and passed \emph{one} adversarial review pass returning no mathematics defect. This is
a strictly weaker tier than \REV{}: one review family, not two. The three results of
\S\ref{sec:adopted} carry it, and the reader should treat them accordingly.

\item \MV{script} --- a purely computational statement: an explicit finite graph together
with its invariants, or a completed bounded sweep. It is evidence of exactly the stated scope
and of nothing beyond it. In particular a \emph{failed} bounded search is never a proof of
impossibility, and \S\ref{sec:limits} records the one occasion on which the campaign
misreported such a search: a sweep over $1112$ local \emph{frames} that found no
counterexample was written up in a way that read as evidence about \emph{graphs}, and a graph
counterexample was found later (Theorem~\ref{thm:g45}).
\end{itemize}

Two further conventions. No result below was obtained with a SAT solver; the only exhaustive
searches are over labelled small-graph tables (explicitly stated), over bounded local
configurations with all edge patterns, and over the neighbourhoods of named explicit graphs.
And where a numerical claim is a \emph{lower bound} on $\pth$ obtained by an early-exit
search, we say ``$\pth \ge k$'' rather than ``$\pth = k$''.

\subsection{Provenance of the arguments}\label{sec:provenance}

The results below were produced by three channels, and it matters which.

The orchestrating agent produced, and is solely responsible for,
Lemmas~\ref{lem:g10}--\ref{lem:g9}, Theorems~\ref{thm:gplus}, \ref{thm:g12},
Lemmas~\ref{lem:g13}--\ref{lem:g34}, Theorems~\ref{thm:g26} and \ref{thm:g35}, the
reductions of \S\ref{sec:peeling}, and every negative result of \S\ref{sec:limits} except
Theorem~\ref{thm:g45}.

The free exploration channel produced first drafts of the extension lemma
Lemma~\ref{lem:g1}, of the radius-$2$ theorem Theorem~\ref{thm:g7}, and of the three results
of \S\ref{sec:adopted}. None was adopted as written: each was re-derived independently before
admission, three claimed solutions of the residual problem of \S\ref{sec:r3} were checked and
\emph{rejected}, and the eventual proof of Theorem~\ref{thm:g35} runs by a route
(Lemma~\ref{lem:g32}) that the free channel never produced. Two garbled justifications inside
the arguments of \S\ref{sec:adopted} were replaced by correct ones during re-derivation and
are marked in place.

The adversarial reviewers produced no positive result by design, with one exception: the
reviewer assigned to \S\ref{sec:adopted} constructed the graph $H$ of
Theorem~\ref{thm:g45}, which refutes a target that the orchestrating agent's own sweep had
recorded as open. That is recorded here as a result of the review channel.

Section~\ref{sec:pocket2} splits the same way. Lemmas~\ref{lem:g41}, \ref{lem:g43},
Theorem~\ref{thm:g42} and Corollary~\ref{cor:g44} are the orchestrating agent's, and they
are the only results in this paper produced by that agent which have been through
\emph{two} adversarial passes from disjoint reviewer families. The witness $Q$ of
Theorem~\ref{thm:g48} came from the free exploration channel, which was asked for the
residual configuration and returned instead a refutation of the reformulation; its recipe was
re-derived and every invariant recomputed here before admission, and its accompanying general
claims were not adopted. The control $W_2$ of Proposition~\ref{prop:w2} is the orchestrating
agent's.

\subsection{Related routes and what is known to be impossible}

Section~\ref{sec:limits} lists five named obstacles, of which one deserves mention up front.
Every bound proved in \S\ref{sec:closed} is driven by $\delta$ or by $\rad$, and both are
destroyed by a single local modification while $\ell$ moves by $O(1/n)$: Hoffman--Singleton
with one edge subdivided is $C_4$-free with $n = 51$, $\delta = 2$, $\rad = 3$,
$\ell = 6.902$, so that $\max(\rad + 2, \delta + 2, \rad + \lfloor \delta/2\rfloor) = 5$
while the conjecture demands $9$. Hence Theorems~\ref{thm:D}, \ref{thm:E}, \ref{thm:F}
\emph{cannot} be combined into a proof; a correct proof must use a quantity as robust as
$\ell$ itself. That observation is why the campaign moved to the $a$-value calculus of
\S\ref{sec:calculus}.

\section{Preliminaries}\label{sec:prelim}

\begin{lemma}[local matching]\label{lem:matching}
\PRV{wowii133\_verify.c}
Let $G$ be $C_4$-free. Then for every $v$ the graph $G[N(v)]$ has maximum degree $\le 1$,
i.e.\ it is a disjoint union of $t(v)$ edges and $d(v) - 2t(v)$ isolated vertices.
Consequently
\[
  a(v) = d(v) - t(v), \qquad \lceil d(v)/2\rceil \le a(v) \le d(v), \qquad
  \ell(G) = \frac{2m - 3T}{n}.
\]
\end{lemma}

\begin{proof}
If $u \in N(v)$ had two neighbours $w_1, w_2 \in N(v)$ then $v w_1 u w_2 v$ is a $4$-cycle on
four distinct vertices. So $G[N(v)]$ is a matching plus isolated vertices; an edge inside
$N(v)$ is exactly a triangle through $v$, so there are $t(v)$ of them. A maximum independent
set takes one endpoint of each edge and every isolated vertex, giving
$t(v) + (d(v) - 2t(v)) = d(v)-t(v)$. Summing and using $\sum_v t(v) = 3T$, $\sum_v d(v) = 2m$
gives the last identity.
\end{proof}

So (C133-b) is the statement $\pth(G) \ge \rad(G) + \lfloor (2m-3T)/n \rfloor$; for graphs of
girth $\ge 5$ it is $\pth \ge \rad + \lfloor \bar d\rfloor$ with $\bar d$ the average degree.
Note $\ell \ge 1$ always, so the open region is $\lfloor \ell\rfloor \ge 3$, i.e.\
$2m - 3T \ge 3n$.

The following vocabulary is used throughout. For $v$ in a $C_4$-free graph, the
\emph{components of $N(v)$} always means the components of $G[N(v)]$, which by
Lemma~\ref{lem:matching} are single vertices and single edges; their number is exactly
$a(v)$. A \emph{geodesic} $u_0 \dots u_d$ is a shortest path, hence automatically induced. A
vertex $x \in N(u_0)$ is a \emph{usable side-neighbour} of the end $u_0$ of that geodesic if
$x$ lies outside $u_1$'s component of $N(u_0)$; one exists iff $a(u_0) \ge 2$, and a usable
side-neighbour satisfies $x \not\sim u_1$ by definition. A vertex $h$ with $a(h) \ge 3$ is
called \emph{high}, or a \emph{hub}, and $H := \{v : a(v) \ge 3\}$.

\begin{lemma}[classification of $a = 1$]\label{lem:g2}
\PRV{wowii133\_gplus.c}
Let $G$ be $C_4$-free. Then $a(v) = 1$ if and only if $N(v)$ is a clique of size $\le 2$,
i.e.\ $v$ is a leaf, or $d(v) = 2$ and the two neighbours of $v$ are adjacent. We call the
latter a \emph{triangle-leaf}.
\end{lemma}

\begin{proof}
$a(v)$ is the number of components of $G[N(v)]$, which is a matching plus isolated vertices;
it equals $1$ iff there is exactly one component, i.e.\ one vertex or one edge.
\end{proof}

\section{The closed part of the conjecture}\label{sec:closed}

\begin{theorem}[the $C_4$-containing branch]\label{thm:A}
\PRV{wowii133\_verify.c}
For every connected $G$, $\pth(G) \ge \rad(G) + 1$. In particular (C133-a) holds.
\end{theorem}

\begin{proof}
Let $c$ be a centre, $\ecc(c) = r = \rad$, and let $z$ satisfy $\mathrm{dist}(c,z) = r$. A
shortest $c$--$z$ path $u_0\dots u_r$ is induced, since a chord $u_iu_j$ with $j \ge i+2$
would give a shorter $c$--$z$ walk. It has $r+1$ vertices.
\end{proof}

\begin{theorem}\label{thm:B}
\PRVn
If $\lfloor \ell(G)\rfloor \le 1$ then (C133-b) holds.
\end{theorem}

\begin{proof}
Immediate from Theorem~\ref{thm:A}.
\end{proof}

\begin{theorem}[radius $\le 1$]\label{thm:C}
\PRV{wowii133\_verify.c}
If $G$ is connected, $C_4$-free, $n \ge 2$ and $\rad(G) \le 1$, then $\ell(G) < 2$; hence
$\lfloor \ell(G)\rfloor = 1$ and (C133-b) holds.
\end{theorem}

\begin{proof}
Let $v$ dominate, so $N(v) = V \setminus \{v\}$. By Lemma~\ref{lem:matching} $G - v$ is a
matching: say $a$ edges and $b$ isolated vertices, $n = 1 + 2a + b$. Then $a(v) = a+b$; a
vertex $u$ matched to $u'$ has $N(u) = \{v, u'\}$ with $v \sim u'$, so $a(u) = 1$; an
unmatched $u$ has $N(u) = \{v\}$, $a(u) = 1$. Hence
$\sum_w a(w) = (a+b) + 2a + b = 3a + 2b < 2(1+2a+b) = 2n$.
\end{proof}

\begin{theorem}[radius $\ge 2$]\label{thm:D}
\PRV{wowii133\_verify.c}
If $G$ is connected, $C_4$-free and $\rad(G) \ge 2$, then $\pth(G) \ge \rad(G) + 2$.
\end{theorem}

\begin{proof}
Write $r = \rad(G)$.

\emph{Case 1: $\diam(G) \ge r+1$.} A geodesic between two vertices at distance $\diam$ is
induced and has $\diam + 1 \ge r+2$ vertices.

\emph{Case 2: $\diam(G) = r$.} Since $r \le \ecc(v) \le \diam$ for every $v$, every vertex is
a centre.

\emph{Case 2a: some vertex $c$ has $d(c) \ge 3$.} Pick $z$ with
$\mathrm{dist}(c,z) = \ecc(c) = r \ge 2$ and a geodesic $c = u_0, u_1, \dots, u_r = z$. At
most one vertex of $N(c)\setminus\{u_1\}$ is adjacent to $u_1$: two such vertices $x, x'$
would give the $4$-cycle $c\,x\,u_1\,x'$. As $|N(c)\setminus\{u_1\}| \ge 2$, choose
$x \in N(c)\setminus\{u_1\}$ with $x \not\sim u_1$. Then $x, c, u_1, \dots, u_r$ is an
induced path on $r+2$ vertices: $x \ne u_i$ for all $i$; $x \not\sim u_1$ by choice;
$x\not\sim u_2$, else $c\,x\,u_2\,u_1$ is a $4$-cycle; and $x \not\sim u_i$ for $i \ge 3$
since $\mathrm{dist}(x,u_i) \ge i-1 \ge 2$.

\emph{Case 2b: $\Delta \le 2$.} Then $G$ is a path or a cycle. A path $P_n$ has
$\diam = n-1$ and $\rad = \lceil (n-1)/2\rceil$, so $\diam = \rad$ forces $n \le 2$,
contradicting $\rad \ge 2$. So $G = C_n$; $C_3$ has radius $1$ and $C_4$ is not $C_4$-free,
so $n \ge 5$, and $\pth(C_n) = n-1 \ge \lfloor n/2\rfloor + 2 = \rad + 2$.
\end{proof}

\begin{corollary}\label{cor:D1}
\PRVn
Conjecture~133 holds for every connected graph containing a $4$-cycle, and for every
connected $C_4$-free graph with $\lfloor \ell(G)\rfloor \le 2$.
\end{corollary}

\begin{proof}
Theorem~\ref{thm:A}; and for $C_4$-free $G$: if $\lfloor\ell\rfloor \le 1$ use
Theorem~\ref{thm:B}; if $\lfloor\ell\rfloor = 2$ then $\rad \ge 2$ by Theorem~\ref{thm:C},
and Theorem~\ref{thm:D} applies.
\end{proof}

Theorem~\ref{thm:D} is tight on $C_5$ and $C_6$. Two greedy bounds complete the picture of
what elementary arguments give.

\begin{theorem}[greedy induced path]\label{thm:E}
\PRV{wowii133\_verify.c}
{\rm (E1)} If $G$ is $C_4$-free and triangle-free (girth $\ge 5$) with $\delta \ge 2$, then
$\pth(G) \ge \delta + 2$.
{\rm (E2)} If $G$ is $C_4$-free with $\delta \ge 3$, then $\pth(G) \ge \delta + 1$.
\end{theorem}

\begin{proof}
Grow an induced path $v_1 \dots v_k$ greedily. A vertex $x \in N(v_k)$ fails to extend it iff
$x$ lies on the path or $x \sim v_i$ for some $i \le k-1$. Count failures: $v_{k-1}$ is the
only path vertex in $N(v_k)$; a common neighbour of $v_{k-1}$ and $v_k$ is a triangle-mate,
of which $C_4$-freeness allows at most one, and none if $G$ is triangle-free; $v_{k-1}$ is
already a common neighbour of $v_{k-2}$ and $v_k$, so $C_4$-freeness gives no other; and each
$v_i$ with $i \le k-3$ has at most one common neighbour with $v_k$. Hence an extension exists
as soon as $d(v_k) \ge k$ ($C_4$-free, $k \ge 3$), resp.\ $d(v_k) \ge k-1$ and $\ge 2$
(girth $\ge 5$). With min degree $\delta$ the greedy runs to $k = \delta+1$, resp.\
$k = \delta + 2$.
\end{proof}

E1 is tight for $C_5$ and for Petersen. In particular every regular graph of girth $\ge 5$
and radius $2$ satisfies the conjecture, since there $\ell = \delta$ exactly and
$\pth \ge \delta + 2 = \rad + \ell$; this covers all Moore graphs of diameter $2$.

\begin{theorem}[geodesic $+$ greedy]\label{thm:F}
\PRV{wowii133\_verify.c}
If $G$ is connected $C_4$-free with $\delta \ge 4$, then
$\pth(G) \ge \rad(G) + \lfloor \delta/2\rfloor$.
\end{theorem}

\begin{proof}
Let $c$ be a centre, $r = \rad$, $z$ at distance $r$, $u_0 \dots u_r$ a geodesic; set
$w_0 = z$ and extend greedily $w_1, w_2, \dots$ at the $z$ end. One step changes the BFS
layer from $c$ by at most $1$, so $w_{j-1}$ lies in a layer $\ge r-j+1$, hence
$\mathrm{dist}(u_i, w_{j-1}) \ge (r-j+1) - i$, so only the $\le j+2$ vertices $u_i$ with
$i \ge r-j-1$ can be within distance $2$ of $w_{j-1}$. Each path vertex $y \ne w_{j-1}$
blocks at most one candidate, since $|N(y) \cap N(w_{j-1})| \le 1$, and one further slot is
lost to $w_{j-2} \in N(w_{j-1})$. Hence at most $(j+2) + (j-2) + 1 = 2j+1$ neighbours of
$w_{j-1}$ are unusable and step $j$ succeeds whenever $d(w_{j-1}) \ge 2j+2$. All steps
$j \le \lfloor\delta/2\rfloor - 1$ succeed.
\end{proof}

\subsection{Falsification and coverage}\label{sec:falsification}

\begin{proposition}\label{prop:falsification}
\MV{wowii133\_verify.c, wowii133\_families.py, wowii133\_sa.c}
No counterexample to Conjecture~133 exists in any of the following ranges:
\begin{center}
\small
\begin{tabular}{@{}llr@{}}
\toprule
range & evaluations & violations \\
\midrule
all connected graphs, $n \le 9$ & $2\,950\,981$ & $0$ \\
all connected graphs, $n \le 10$ & $136\,362\,861$ & $0$ \\
all connected $C_4$-free graphs, $n \le 13$ & $13\,544\,784$ & $0$ \\
$79$ structured graphs, $n \le 183$ & --- & $0$ \\
simulated annealing, $C_4$-free, $n = 10..24$ & $\approx 2\cdot 10^{6}$ proposals & $0$ \\
\bottomrule
\end{tabular}
\end{center}
\end{proposition}

The enumerator generates connected graphs by vertex augmentation, which is complete because
every connected graph on $k+1 \ge 2$ vertices has a non-cut vertex; isomorph rejection uses a
canonical certificate, and at the last level no rejection is performed, so that level is a
superset of all classes. The generator self-validates: its class counts
$1, 2, 6, 21, 112, 853, 11117, 261080$ for $n = 2..9$ are the numbers of connected graphs on
$n$ nodes, and in $C_4$-free mode the full enumerator and the $C_4$-free-only enumerator agree
exactly ($36\,810$ graphs up to $n = 10$). The structured family includes Petersen, the
Hoffman--Singleton graph, the incidence graphs of $PG(2,q)$ for $q \le 9$, the
Erd\H{o}s--R\'enyi orthogonal polarity graphs $ER_q$ for $q \le 13$ --- the extremal
$C_4$-free graphs, and the most dangerous family since they maximise $\ell$ at radius $2$ ---
cages and random girth-$5$ regular graphs.

\begin{proposition}[equality cases]\label{prop:tight}
\MV{wowii133\_verify.c}
Over the whole exhaustive range the $C_4$-free branch is tight for exactly five graphs:
\begin{center}
\small
\begin{tabular}{@{}lrrrrrr@{}}
\toprule
$G$ & $n$ & $\rad$ & $\ell$ & $\lfloor\ell\rfloor$ & $\pth$ & RHS \\
\midrule
$K_2$ & $2$ & $1$ & $1$ & $1$ & $2$ & $2$\\
$K_3$ & $3$ & $1$ & $1$ & $1$ & $2$ & $2$\\
$C_5$ & $5$ & $2$ & $2$ & $2$ & $4$ & $4$\\
$C_6$ & $6$ & $3$ & $2$ & $2$ & $5$ & $5$\\
Petersen & $10$ & $2$ & $3$ & $3$ & $5$ & $5$\\
\bottomrule
\end{tabular}
\end{center}
Consequently no argument that loses even one vertex on these five can prove the conjecture.
\end{proposition}

Measured against Theorems~\ref{thm:A}, \ref{thm:D}, \ref{thm:E}, \ref{thm:F} the coverage of
the exhaustive $C_4$-free range is $100\%$ for $n \le 10$ and $99.9955\%$ for $n \le 13$,
leaving $606$ graphs ($516$ with $\rad = 2$, $90$ with $\rad = 3$, all with
$\lfloor\ell\rfloor = 3$), each of them one vertex short. There are $622$ $C_4$-free graphs
with $\lfloor \ell\rfloor = 3$ on $n \le 13$ vertices at all, and none with
$\lfloor\ell\rfloor \ge 4$ below $n = 14$. Theorem~\ref{thm:gplus} below closes all $606$:
all $622$ have $\mu \ge 2$, and the residual over the whole exhaustive $C_4$-free range
becomes $0$.

\subsection{Auxiliary claims: verified, not proved}\label{sec:aux}

The same machinery tested four statements that are \emph{not} used anywhere in this paper and
are recorded because one of them is the natural target.

\begin{proposition}\label{prop:aux}
\MV{wowii133\_verify.c}
The following hold over the ranges stated, with the number of violations shown.
\begin{center}
\small
\begin{tabular}{@{}llll@{}}
\toprule
& claim & range verified & violations \\
\midrule
A1 & every connected $G$: $\pth \ge 2\rad - 1$ & all connected $n \le 10$ & $0$\\
A4 & $C_4$-free: $\pth \ge \rad + \delta - 1$ & all $C_4$-free $n \le 13$ & $0$\\
A6 & $C_4$-free, $G \ne K_3$: $\pth \ge \rad + \delta$ & all $C_4$-free $n \le 13$ &
   $1$ ($K_3$)\\
A7 & $C_4$-free: $\ell(G) \le \pth(G) - \rad(G)$ & all $C_4$-free $n \le 13$;
   all connected $n \le 10$ & $0$\\
\bottomrule
\end{tabular}
\end{center}
None of these is proved. A7 is the floor-free strengthening of (C133-b), i.e.\
$2m - 3T \le n(\pth - \rad)$; it is tight on the same five graphs as
Proposition~\ref{prop:tight} and is false without $C_4$-freeness ($K_{3,3}$ has $\ell = 3$
and $\pth - \rad = 1$).
\end{proposition}

Theorem~\ref{thm:g12} below proves a non-trivial instance of A7.

\section{An endpoint-extension calculus}\label{sec:calculus}

All lemmas in this section are unconditional facts about $C_4$-free graphs. They are the
tools that replace the degree-driven bounds of \S\ref{sec:closed}, which the surgery obstacle
of \S\ref{sec:limits} rules out.

\begin{lemma}[one-vertex extension]\label{lem:g10}
\PRV{wowii133\_gplus.c}
Let $G$ be $C_4$-free with a geodesic $u_0 \dots u_d$, $d \ge 2$, and $a(u_d) \ge 2$. Then
$\pth(G) \ge d+2$.
\end{lemma}

\begin{proof}
By Lemma~\ref{lem:matching} the number of components of $N(u_d)$ is $a(u_d) \ge 2$; take $x$
in a component other than $u_{d-1}$'s. Then $x \ne u_i$ for all $i$, since
$\mathrm{dist}(u_d,x) = 1$ while $\mathrm{dist}(u_d,u_i) = d-i$, and $i = d-1$ is excluded;
$x \not\sim u_{d-1}$ (different components); $x \not\sim u_{d-2}$, else
$x, u_d, u_{d-1}, u_{d-2}$ are four distinct vertices forming a $4$-cycle; and for
$i \le d-3$, $\mathrm{dist}(x,u_i) \ge (d-i)-1 \ge 2$. So $u_0 \dots u_d\,x$ is induced.
\end{proof}

\begin{lemma}[two-vertex extension]\label{lem:g1}
\PRV{wowii133\_gplus.c}
Let $G$ be $C_4$-free with a geodesic $u_0\dots u_d$, $d \ge 2$, $a(u_0) \ge 3$ and
$a(u_d) \ge 2$. Then $\pth(G) \ge d+3$.
\end{lemma}

\begin{proof}
Take $x_1, x_2$ in two components of $N(u_0)$ other than $u_1$'s (possible since
$a(u_0) \ge 3$) and $y$ in a component of $N(u_d)$ other than $u_{d-1}$'s. If $y$ were
adjacent to both $x_1$ and $x_2$ then $u_0$ and $y$ would have two common neighbours, a
$C_4$; so fix $x \in \{x_1,x_2\}$ with $x \not\sim y$.

We check that $x, u_0, u_1, \dots, u_d, y$ is induced. As in Lemma~\ref{lem:g10},
$y \not\sim u_{d-1}$, $y\not\sim u_{d-2}$ and $y \not\sim u_i$ for $i \le d-3$, and
symmetrically $x \not\sim u_1$, $x \not\sim u_2$ (else $u_0, u_2$ have the two common
neighbours $x, u_1$), and $x \not\sim u_i$ for $i \ge 3$ by distance. For $d = 2$ the only
remaining pair is $(y,u_0)$: $y \sim u_0$ would give $u_0, u_2$ the two common neighbours
$u_1, y$. For $d \ge 3$, $\mathrm{dist}(u_0,y) \ge d - 1 \ge 2$. Finally $x \ne y$ and
$x \not\sim y$ by choice, and $x, y \notin \{u_0,\dots,u_d\}$ by the distance arguments
above.
\end{proof}

\begin{lemma}[both ends, $d \ge 4$]\label{lem:g9}
\PRV{wowii133\_gplus.c}
Let $G$ be $C_4$-free with a geodesic $u_0 \dots u_d$, $d \ge 4$, $a(u_0)\ge 2$ and
$a(u_d) \ge 2$. Then $\pth(G) \ge d+3$. The hypothesis $d \ge 4$ cannot be dropped.
\end{lemma}

\begin{proof}
Take $x \in N(u_0)$ outside $u_1$'s component and $y \in N(u_d)$ outside $u_{d-1}$'s
component. Lemma~\ref{lem:g10}'s checks, applied at both ends, give $x,y \notin
\{u_0,\dots,u_d\}$ and non-adjacency of $x$ and of $y$ to every $u_i$ they must avoid.
Moreover $\mathrm{dist}(x,y) \ge \mathrm{dist}(u_0,u_d) - 2 = d - 2 \ge 2$, so $x \ne y$ and
$x \not\sim y$. Hence $x, u_0, \dots, u_d, y$ is induced.
\end{proof}

For $d = 3$ the statement is false: over the exhaustive $C_4$-free range $n \le 13$ there are
$334$ pairs $(u,z)$ at distance $3$ with both $a$-values $\ge 2$ and $\pth < 6$; the two
extensions may be adjacent, producing an induced $C_6$ rather than a $P_6$. \MV{wowii133\_gplus.c}

\begin{theorem}\label{thm:g}
\PRV{wowii133\_gplus.c}
Every connected $C_4$-free graph satisfies $\pth(G) \ge \rad(G) + \min\{3, \mu(G)\}$.
\end{theorem}

\begin{proof}
If $\rad \le 1$ then $\ell < 2$ by Theorem~\ref{thm:C}, hence $\mu \le \ell < 2$ so
$\mu = 1$, and Theorem~\ref{thm:A} suffices. Let $\rad \ge 2$. If $\mu = 1$ use
Theorem~\ref{thm:A}; if $\mu = 2$ use Theorem~\ref{thm:D}; if $\mu \ge 3$, take a centre $c$
and $z$ at distance $\rad$ and apply Lemma~\ref{lem:g1} to a $c$--$z$ geodesic.
\end{proof}

Theorem~\ref{thm:g} is tight on Petersen. Two immediate corollaries: a triangle-free
$C_4$-free graph with $\delta \ge 3$ has $\pth \ge \rad + 3$ (there $a = d \ge 3$), and so
does any $C_4$-free graph with $\delta \ge 5$ (there $a \ge \lceil d/2\rceil \ge 3$).

The next theorem is the workhorse: it removes the gap at $\mu = 2$, where
Theorem~\ref{thm:g} only gives $\rad + 2$.

\begin{theorem}\label{thm:gplus}
\PRV{wowii133\_gplus.c}
Let $G$ be connected and $C_4$-free with $\mu(G) \ge 2$ and with some vertex $h$ satisfying
$a(h) \ge 3$. Then $\pth(G) \ge \ecc(h) + 3 \ge \rad(G) + 3$, provided $\rad(G) \ge 2$.
\end{theorem}

\begin{proof}
Take $z$ with $\mathrm{dist}(h,z) = \ecc(h) =: d \ge \rad \ge 2$. Then $a(h) \ge 3$ and
$a(z) \ge \mu \ge 2$, so Lemma~\ref{lem:g1} applies to an $h$--$z$ geodesic.
\end{proof}

\begin{corollary}\label{cor:gplus1}
\PRV{wowii133\_gplus.c}
Let $G$ be connected and $C_4$-free with $\ell(G) \ge 3$ and no vertex of $a$-value $1$
(equivalently, by Lemma~\ref{lem:g2}, no leaf and no triangle-leaf). Then
$\pth(G) \ge \rad(G) + 3$. In particular (C133-b) holds for all such graphs with
$\lfloor \ell\rfloor \le 3$.
\end{corollary}

\begin{proof}
$\ell \ge 3$ gives $\sum_v a(v) \ge 3n$, so some $a(h) \ge 3$; and $\ell \ge 3 > 2$ with
Theorem~\ref{thm:C} gives $\rad \ge 2$. Apply Theorem~\ref{thm:gplus}.
\end{proof}

This is what closes the $606$ residual graphs of \S\ref{sec:falsification}: all $622$
$C_4$-free graphs with $\lfloor\ell\rfloor = 3$ on at most $13$ vertices have $\mu \ge 2$.

Finally, a lemma that reaches one vertex further at the cost of a hypothesis on a
side-neighbour. It is the seed of \S\ref{sec:adopted}.

\begin{lemma}[deep endpoint extension, $d \ge 5$]\label{lem:g14}
\PRV{w133\_r3\_counterexample.py}
Let $G$ be $C_4$-free with a geodesic $u_0 \dots u_d$, $d \ge 5$, $a(u_0) \ge 3$,
$a(u_d) \ge 2$, and a usable side-neighbour $x$ of $u_0$ with $a(x) \ge 4$. Then
$\pth(G) \ge d + 4$. In particular $\mu \ge 4$ and $\rad \ge 5$ imply $\pth \ge \rad + 4$.
\end{lemma}

\begin{proof}
Choose $x$ as in the hypothesis and $y \in N(u_d)$ outside $u_{d-1}$'s component. The graph
$N(x)$ has $\ge a(x) \ge 4$ components; $u_0$'s is one. Since $x \not\sim u_2$ (else
$u_0, u_2$ have common neighbours $x, u_1$) and $x \not\sim u_3$ (distance),
$C_4$-freeness gives $|N(x) \cap N(u_2)| \le 1$ and $|N(x) \cap N(u_3)| \le 1$, so at most
one further component is destroyed by each; take $w$ in a surviving component. Then
$w \not\sim u_0$ (different components of $N(x)$); $w \not\sim u_1$, else $w, x, u_0, u_1$
form a $C_4$; $w \not\sim u_2, u_3$ by the choice of component; and
$\mathrm{dist}(w,u_i) \ge i - 2 \ge 2$ for $i \ge 4$. Also $w \notin \{u_0,\dots,u_d\}$
since $\mathrm{dist}(u_0,w) = 2$ and $w \ne u_2$ because $x \not\sim u_2$. Finally
$\mathrm{dist}(w,y) \ge (d-1) - 2 = d-3 \ge 2$ and $\mathrm{dist}(x,y) \ge (d-1)-1 = d-2 \ge
3$, both using $d \ge 5$. So $w, x, u_0, \dots, u_d, y$ is induced on $d+4$ vertices.
\end{proof}

The two steps marked ``$d \ge 5$'' are exactly what fails at $d = 4$ and $d = 3$;
\S\ref{sec:adopted} repairs the first and \S\ref{sec:limits} shows the second cannot be
repaired for free.

\section{The radius-2 layer}\label{sec:r2}

Write $(\star)$ for the assertion that $G$ has \emph{no} pair of vertices at distance exactly
$\rad(G)$ whose $a$-values are $(\ge 3, \ge 2)$ in some order. By Lemma~\ref{lem:g1},
$(\star)$ is a necessary condition on any graph with $\ell \ge 3$ and
$\pth < \rad + 3$, and it is the hypothesis under which the following two lemmas run.

\begin{lemma}\label{lem:g4}
\PRV{w133\_res\_kill.py}
Assume $(\star)$, and let $r = \rad(G)$ in a connected $C_4$-free graph. If $a(h) \ge 3$ then
$\ecc(h) = r$.
\end{lemma}

\begin{proof}
Suppose $\ecc(h) \ge r+1$ and take a geodesic $p_0 = h, p_1, \dots, p_e$ with
$e = \ecc(h)$. The vertex $p_r$ is at distance exactly $r$ from $h$, so $(\star)$ forces
$a(p_r) \le 1$, hence $a(p_r) = 1$; by Lemma~\ref{lem:g2}, $N(p_r)$ is a clique, and since it
contains $p_{r-1}$ and $p_{r+1}$ we get $p_{r-1} \sim p_{r+1}$. Then
$h \to p_{r-1} \to p_{r+1}$ is a walk of length $r$, contradicting
$\mathrm{dist}(h,p_{r+1}) = r+1$.
\end{proof}

\begin{lemma}\label{lem:g5}
\PRV{w133\_res\_kill.py}
Assume $(\star)$. If $a(h) \ge 3$ and $h = p_0, \dots, p_r = z$ is a geodesic to a vertex at
distance $r = \rad$, then $a(p_i) \ge 2$ for every interior vertex $1 \le i \le r-1$.
\end{lemma}

\begin{proof}
$a(p_i) = 1$ would make $N(p_i)$ a clique containing $p_{i-1}$ and $p_{i+1}$, hence
$p_{i-1} \sim p_{i+1}$ and the geodesic could be shortened.
\end{proof}

So in any graph satisfying $(\star)$ a geodesic from a high vertex to the periphery has the
rigid $a$-profile $(\ge 3), (\ge 2), \dots, (\ge 2), (1)$. The radius-$2$ case of that
rigidity is contradictory.

\begin{theorem}[the radius-$2$ layer]\label{thm:g7}
\PRV{w133\_res\_kill.py}
Let $G$ be connected and $C_4$-free with $\rad(G) = 2$ and $\ell(G) \ge 3$. Then $G$ has a
pair of vertices at distance exactly $2$ whose $a$-values are $(\ge 3, \ge 2)$.
\end{theorem}

\begin{proof}
Suppose not; then $(\star)$ holds with $r = 2$. From $\ell \ge 3$ we get
$\sum_v a(v) \ge 3n$, so some vertex $h$ has $a(h) \ge 3$; by Lemma~\ref{lem:g4},
$\ecc(h) = 2$. Put $A := N(h)$, $m := |A|$, and $B := V \setminus (\{h\} \cup A)$, the
distance-$2$ layer. By $(\star)$, every $b \in B$ has $a(b) \le 1$, hence $a(b) = 1$.

\emph{Step 1: unique parents.} Every $b \in B$ has at least one neighbour in $A$ (as
$\mathrm{dist}(h,b) = 2$) and at most one, since two would give $h$ and $b$ two common
neighbours. Write $p(b) \in A$ for that parent and $B_\alpha := p^{-1}(\alpha)$ for the fibre
of $\alpha \in A$.

\emph{Step 2: no edges between distinct fibres.} If $b \in B_\alpha$, $c \in B_\beta$ with
$\alpha \ne \beta$ and $b \sim c$, then $N(b)$ contains $\alpha$ and $c$, which are
non-adjacent (the unique parent of $c$ is $\beta \ne \alpha$), so $a(b) \ge 2$, contradicting
$a(b) = 1$.

\emph{Step 3: the two matchings.} $G[A] \subseteq G[N(h)]$ and
$G[B_\alpha] \subseteq G[N(\alpha)]$ are matchings plus isolated vertices, by
Lemma~\ref{lem:matching}. Let $p$ be the number of edges of $G[A]$, and for
$\alpha \in A$ let $s_\alpha$ and $q_\alpha$ be the numbers of isolated vertices and of edges
of $G[B_\alpha]$; put $S := \sum_\alpha s_\alpha$ and $Q := \sum_\alpha q_\alpha$.

\emph{Step 4: the count.} First $a(h) = \comp(A) = m - p$. Next fix $\alpha \in A$. By
Step~2, $N(\alpha) = \{h\} \cup M(\alpha) \cup B_\alpha$, where $M(\alpha)$ is $\alpha$'s
partner in $G[A]$ if it has one; $h$ is adjacent to $M(\alpha)$ and to nothing in
$B_\alpha$, and no vertex of $B_\alpha$ is adjacent to $M(\alpha)$ (it would then have two
parents). Hence the components of $N(\alpha)$ are $\{h\} \cup M(\alpha)$ together with the
$s_\alpha + q_\alpha$ components of $G[B_\alpha]$, so
$a(\alpha) = 1 + s_\alpha + q_\alpha$. Finally $a(b) = 1$ for $b \in B$ and
$|B| = S + 2Q$. Therefore
\[
  \sum_v a(v) = (m-p) + (m + S + Q) + (S+2Q) = 2m - p + 2S + 3Q,
  \qquad n = 1 + m + S + 2Q,
\]
so that
\[
  3n - \sum_v a(v) = 3 + m + p + S + 3Q \;\ge\; 3 \;>\; 0,
\]
i.e.\ $\ell(G) < 3$, contradicting the hypothesis.
\end{proof}

\begin{remark}
The proof never uses any hypothesis on $\pth(G)$: the contradiction comes from $(\star)$,
$\rad = 2$ and the existence of a high vertex alone. This is what makes
Theorem~\ref{thm:g7} composable with Lemma~\ref{lem:g1}.
\end{remark}

\begin{corollary}[the radius-$2$ layer is closed]\label{cor:g8}
\PRV{w133\_res\_kill.py}
Let $G$ be connected and $C_4$-free with $\rad(G) = 2$ and $\ell(G) \ge 3$. Then
$\pth(G) \ge 5 = \rad(G) + 3$. In particular (C133-b) holds for every connected $C_4$-free
graph with $\rad = 2$ and $\lfloor \ell\rfloor = 3$.
\end{corollary}

\begin{proof}
Theorem~\ref{thm:g7} supplies a pair at distance $2$ with $a$-values $(\ge 3,\ge 2)$; apply
Lemma~\ref{lem:g1} with $d = 2$.
\end{proof}

\section{Peeling induction and the residual problem}\label{sec:peeling}

From here on the target is
\begin{equation}\label{eq:X}
  \textbf{(X)}\qquad G \text{ connected } C_4\text{-free},\; \ell(G) \ge 3
  \;\Longrightarrow\; \pth(G) \ge \rad(G) + 3,
\end{equation}
which implies (C133-b) for $\lfloor \ell\rfloor \le 3$ and is the whole content of the band
$\lfloor\ell\rfloor = 3$. Corollary~\ref{cor:gplus1} settles \eqref{eq:X} when there is no
vertex of $a$-value $1$; the induction below reduces the general case to a residual that can
be attacked.

\begin{lemma}[deleting an $a = 1$ vertex]\label{lem:g6}
\PRV{wowii133\_gplus.c}
Let $G$ be $C_4$-free and $a(v) = 1$. Then $G - v$ is connected, $C_4$-free, and preserves
all distances between the remaining vertices.
\end{lemma}

\begin{proof}
By Lemma~\ref{lem:g2}, $v$ is a leaf or a triangle-leaf. A leaf is never an interior vertex
of a shortest path; a triangle-leaf $v$ with $N(v) = \{x,y\}$, $x \sim y$, can always be
bypassed by the edge $xy$. So no shortest path between surviving vertices needs $v$.
\end{proof}

\begin{lemma}[peeling arithmetic]\label{lem:g11}
\PRV{w133\_r11\_adjudicate.py}
Let $G$ be $C_4$-free with $a(v) = 1$, and write $\Sigma(G) := \sum_u a(u)$. Then
\[
  \Sigma(G-v) - 3(n-1) \;=\; \bigl(\Sigma(G) - 3n\bigr) + c, \qquad
  c = \begin{cases} 1 & v \text{ a leaf},\\ 2 & v \text{ a triangle-leaf}.\end{cases}
\]
In particular the slack $\Sigma - 3n$ is strictly increasing under peeling, by at least $1$
per deleted vertex.
\end{lemma}

\begin{proof}
If $v$ is a leaf with neighbour $x$, then $v$ is an isolated vertex of $G[N(x)]$, so $a(x)$
drops by exactly $1$ and no other $a$-value changes: $\Sigma$ drops by $2$ and $n$ by $1$,
so the slack rises by $1$. If $v$ is a triangle-leaf with $N(v) = \{x,y\}$, $x \sim y$, then
$\{v,y\}$ is an edge component of $G[N(x)]$ which degenerates to the isolated vertex $y$;
its contribution to $a(x)$ is still $1$, so $a(x)$ is unchanged, and symmetrically $a(y)$ is
unchanged. Hence $\Sigma$ drops by $1$ and $n$ by $1$, and the slack rises by $2$.
\end{proof}

\begin{remark}[a corrected statement]\label{rem:b10}
Earlier internal versions of the project stated Lemma~\ref{lem:g11} as ``the slack rises by
exactly $1$''. That is \emph{false} for triangle-leaves, as the adversarial review of
\S\ref{sec:adopted} pointed out and as the explicit graph CE-1 of \S\ref{sec:limits} shows:
deleting its leaf $1$ moves $(\Sigma,n)$ from $(19,10)$ to $(17,9)$, slack $-11 \to -10$;
deleting its triangle-leaf $8$ moves it to $(18,9)$, slack $-11 \to -9$. Every use in this
paper needs only the monotone direction ``$\ge +1$'', so no conclusion moves; but the
``exactly'' was wrong and is corrected here.
\end{remark}

\begin{proposition}[the induction skeleton]\label{prop:skeleton}
\PRVn
Let $G$ be a counterexample to \eqref{eq:X} of minimum order, and put
$\Lambda := \{v : a(v) = 1\}$. Then:
\begin{enumerate}
\item[(i)] $\rad(G) \ge 2$ and $\Lambda \ne \emptyset$;
\item[(ii)] deleting any $v \in \Lambda$ strictly lowers the radius;
\item[(iii)] the \emph{core} $C := G - \Lambda$ is connected, $C_4$-free, isometrically
embedded in $G$, every vertex of $C$ has $a_G \ge 2$, and $\ell(C) > 3$;
\item[(iv)] $\diam(G) \le \rad(G) + 1$, and $(\star)$ holds in $G$.
\end{enumerate}
\end{proposition}

\begin{proof}
(i) $\ell \ge 3 > 2$ and Theorem~\ref{thm:C} give $\rad \ge 2$; $\Lambda = \emptyset$ would
contradict Corollary~\ref{cor:gplus1}. (ii) If $\rad(G-v) = \rad(G)$ for some $v \in \Lambda$
then $G - v$ is connected, $C_4$-free (Lemma~\ref{lem:g6}) with $\ell \ge 3$
(Lemma~\ref{lem:g11}) and smaller order, so by minimality
$\pth(G-v) \ge \rad(G-v) + 3 = \rad(G)+3$, and $\pth(G) \ge \pth(G-v)$, a contradiction.
(iii) Connectivity and isometry are Lemma~\ref{lem:g6} applied to the whole of $\Lambda$ at
once, whose proof shows directly that no shortest path between core vertices uses a vertex of
$\Lambda$. Deleting the members of $\Lambda$ one at a time --- each remains of $a$-value $1$
when its turn comes, a triangle-leaf possibly having become a leaf --- and applying
Lemma~\ref{lem:g11} $|\Lambda| \ge 1$ times gives
$\Sigma(C) - 3n_C \ge (\Sigma(G) - 3n) + |\Lambda| \ge 1$, i.e.\ $\ell(C) > 3$.
(iv) If $\diam \ge \rad + 2$ a diametral geodesic is an induced path on
$\diam + 1 \ge \rad + 3$ vertices. $(\star)$ is Lemma~\ref{lem:g1} together with
$\pth(G) \le \rad(G)+2$.
\end{proof}

For $\rad \ge 4$ one obtains in addition, from Lemma~\ref{lem:g9}, that all vertices of
$a$-value $\ge 2$ lie pairwise within distance $\rad - 1$, whence $\diam(C) \le \rad - 1$ and
$C$ is self-centred with $\rad(C) = \diam(C) = \rad(G) - 1 =: D$; moreover every core vertex
lies at distance exactly $D$ from some vertex of the ``parent set''
$P := \{$core vertices with a $\Lambda$-neighbour$\}$. Applying minimality to $C$ gives
$\pth(C) \ge D + 3 = \rad(G)+2$, one vertex short. This is the residual:

\begin{problem}[the residual problem]\label{prob:rp}
Let $C$ be connected, $C_4$-free, self-centred with $\rad(C) = \diam(C) = D$ and
$\ell(C) > 3$, and let $P \subseteq V(C)$ satisfy: for every $u$ there is $p \in P$ with
$\mathrm{dist}(u,p) = D$. Show that $C$ has an induced path on $D+3$ vertices with at least
one end in $P$ (or one on $D+2$ vertices with both ends in $P$).
\end{problem}

\begin{problem}[the strengthened form]\label{prob:rpd}
$C$ connected $C_4$-free, $\diam(C) = D$, $\ell(C) > 3$ $\Longrightarrow$
$\pth(C) \ge D + 4$.
\end{problem}

Problem~\ref{prob:rpd} implies Problem~\ref{prob:rp}, since $D+4 > D+3$ makes the endpoint
condition free, and it is the integer form of the strengthening A7 of
Proposition~\ref{prop:aux} in the self-centred case. Its case $D = 2$ is the following
theorem; $D \ge 5$ is available from Lemma~\ref{lem:g14} under an extra side-neighbour
hypothesis, $D = 4$ from Lemma~\ref{lem:g36} below, and $D = 3$ is genuinely open --- indeed
\S\ref{sec:limits} shows that the local route to it is closed.

\begin{theorem}[$\diam = 2$]\label{thm:g12}
\PRV{wowii133\_gplus.c, w133\_g12\_nonempty.py}
Let $C$ be connected, $C_4$-free, with $\diam(C) = 2$ and $\ell(C) > 3$ (strictly). Then
$\pth(C) \ge 6$. Since $\diam = 2$ and $\ell > 2$ force $\rad(C) = 2$ by
Theorem~\ref{thm:C}, this says $\pth \ge \rad + 4$.
\end{theorem}

\begin{proof}
\emph{Step 0.} $\diam = 2$ and $C_4$-freeness make any two non-adjacent vertices have
\emph{exactly} one common neighbour.

\emph{Step 1: the triangle-free case.} We first record the step that an adversarial review
found missing in an earlier version of this proof.

\begin{quote}
\emph{Claim.} If $C$ is connected, $C_4$-free, triangle-free with $\diam(C) = 2$ and
$\ell(C) > 3$, then $C$ contains a cycle and its girth is exactly $5$.\\
\emph{Proof.} If $C$ were acyclic it would be a tree of diameter $2$, i.e.\ a star
$K_{1,m}$, and $\ell(K_{1,m}) = 2m/(m+1) < 2$, contradicting $\ell > 3$. A graph containing a
cycle satisfies $\gir \le 2\diam + 1 = 5$, and triangle-freeness with $C_4$-freeness gives
$\gir \ge 5$.
\end{quote}

\noindent Without this step the claim ``girth $5$'' is false: a star is triangle-free,
$C_4$-free and of diameter $2$ but has no cycle at all, and is not regular. Given girth $5$
and diameter $2$, $C$ is a Moore graph, hence $k$-regular with $k \in \{2,3,7,57\}$; being
triangle-free, $\ell(C) = k$, so $\ell > 3$ forces $k \in \{7,57\}$ and $\delta \ge 7$, and
Theorem~\ref{thm:E}(E1) gives $\pth \ge \delta + 2 \ge 9 \ge 6$.

\emph{Step 2: the case with a triangle.} Suppose $\pth(C) \le 5$ and let $\{a,b,c\}$ be a
triangle. Put $A := N(a)\setminus\{b,c\}$, $B := N(b)\setminus\{a,c\}$,
$C_s := N(c)\setminus\{a,b\}$ and $R := V \setminus (\{a,b,c\}\cup A \cup B \cup C_s)$.

(i) \emph{$A, B, C_s$ are pairwise disjoint and pairwise non-adjacent.} $x \in A \cap B$
would make $c, x$ two common neighbours of $a,b$; and $\alpha \in A$, $\beta \in B$ with
$\alpha \sim \beta$ would make $b, \alpha$ two common neighbours of $a, \beta$.

(ii) \emph{Every $u \in R$ has exactly one neighbour in each of $A$, $B$, $C_s$.} As
$u \not\sim a$ and $\diam = 2$, $u$ and $a$ have exactly one common neighbour
$w \in \{b,c\} \cup A$; since $u \not\sim b,c$ we get $w \in A$.

(iii) \emph{Counting identities.} $a(a) = 1 + \comp(A)$, and for $\alpha \in A$,
$a(\alpha) = 1 + \comp(R_\alpha)$ where $R_\alpha := N(\alpha)\cap R$. Indeed
$N(a) = \{b,c\}\cup A$ whose internal edges are $bc$ together with the edges of $A$ (no
$b$--$A$ edges by (i)); and $N(\alpha) = \{a\} \cup M(\alpha) \cup R_\alpha$ with internal
edges $a$--$M(\alpha)$ and those of $R_\alpha$, since $M(\alpha) \subseteq A$ has no
$R_\alpha$-neighbour by (ii).

(iv) \emph{$R$ is a clique.} Let $\rho,\sigma \in R$ be non-adjacent, and let $f_A, f_B,
f_C$ denote the unique neighbours supplied by (ii). If $f_A(\rho)\ne f_A(\sigma)$ and
$f_B(\rho) \ne f_B(\sigma)$, put $\alpha := f_A(\rho)$, $\beta := f_B(\sigma)$; then
$\rho, \alpha, a, b, \beta, \sigma$ is an induced $P_6$ --- $\rho \not\sim a,b$ by the
definition of $R$; $\rho \not\sim\beta$ since $\rho$'s unique $B$-neighbour differs from
$\beta$; $\alpha \not\sim b, \beta, \sigma$ and $a \not\sim \beta,\sigma$, $b \not\sim\sigma$
by (i),(ii) --- contradicting $\pth \le 5$. Using the edges $ac$ and $bc$ likewise, at least
two of the three coordinates coincide; but two coincidences give $\rho,\sigma$ two common
neighbours, a $C_4$.

(v) $|R| \le 3$, since a clique on $4$ vertices contains a $C_4$.

(vi) \emph{The count.} Write $s := |A|+|B|+|C_s|$ and $\rho := |R|$. By (iii),
$\sum_v a(v) = 3 + \sum_{T}\comp + s + \sum_\alpha \comp(R_\alpha) + \sum_{u \in R} a(u)$,
where $\sum_T \comp \le s$ and $\sum_\alpha \comp(R_\alpha) \le 3\rho$ (each $R$-vertex lies
in exactly one fibre on each of the three sides). Now $n = 3 + s + \rho$, and we compare with
$3n$:
\begin{itemize}
\item $\rho = 0$: $a(\alpha) = 1$ for all $\alpha$, so $\sum a \le 3 + 2s$ while
$3n = 9+3s$.
\item $\rho = 1$: $d(u) = 3$ so $a(u) \le 3$, $\sum a \le 9 + 2s$ while $3n = 12+3s$.
\item $\rho = 2$: $d(u) = 4$ so $a(u) \le 4$, $\sum a \le 17+2s$ while $3n = 15+3s$, forcing
$s < 2$; but (ii) gives $s \ge 3$.
\item $\rho = 3$: $R$ is a triangle, and for $u \in R$,
$N(u) = \{\alpha_u,\beta_u,\gamma_u\} \cup (R \setminus\{u\})$ where $R\setminus\{u\}$ is an
edge; by the matching property $\alpha_u,\beta_u,\gamma_u$ are non-adjacent to
$R\setminus\{u\}$, so $a(u) = 4$ exactly, and $f_A$ is injective on $R$, giving
$|A|,|B|,|C_s| \ge 3$ and $s \ge 9$. Then $\sum a \le 3+2s+9+12 = 24+2s$ while
$3n = 18+3s$, forcing $s < 6$.
\end{itemize}
In every case $\sum_v a(v) \le 3n$, contradicting $\ell(C) > 3$.
\end{proof}

\begin{remark}[status of Theorem~\ref{thm:g12}]
Step 1's Claim was added after an adversarial review pass found the gap; the star
counterexample above is that reviewer's. The patched statement has \emph{not} itself been
through a second, cross-family review pass --- that item is still queued in the campaign's
ledger --- which is why Theorem~\ref{thm:g12} carries \PRV{} and not \REV{}. The structural
steps (i)--(iv) were asserted over the exhaustive $C_4$-free range $n \le 13$ on $422$
non-vacuous diameter-$2$ instances with $0$ failures; the theorem's own hypothesis class is
empty in that range (the closest graph, Petersen, has $\ell = 3$ exactly and is excluded by
strictness), so its positive evidence is the explicit witness below.
\end{remark}

\begin{proposition}[non-vacuity of Theorem~\ref{thm:g12}]\label{prop:g12nonempty}
\MV{w133\_g12\_nonempty.py, w133\_res3\_search.py}
Both branches of Theorem~\ref{thm:g12}'s hypothesis class are inhabited. The
Hoffman--Singleton graph has $n = 50$, girth $5$, $\diam = 2$, is $7$-regular and
triangle-free, so $\ell = 7 > 3$; it satisfies the conclusion with an induced-path
certificate of $20$ vertices. The Erd\H{o}s--R\'enyi polarity graphs $ER_5, ER_7, ER_{11}$
are $C_4$-free of diameter $2$ \emph{with} triangles and $\ell = 3.871, 4.912, 6.947$
respectively, with induced-path certificates $15, 22, 37$.
\end{proposition}

\section{The radius-3 layer}\label{sec:r3}

This section proves the main theorem of the paper.

\begin{theorem}[the radius-$3$ layer]\label{thm:r3}
\PRVn
Let $G$ be a counterexample to \eqref{eq:X} of minimum order. Then $\rad(G) \ge 4$.
\end{theorem}

The proof occupies \S\ref{sec:reduction}--\S\ref{sec:hexagon}. It reduces the case
$\rad(G) = 3$ to the emptiness of an explicit class, then proves that emptiness in two
theorems.

\subsection{The reduction}\label{sec:reduction}

\begin{definition}
Write (R3b) for the property: \emph{every vertex of $a$-value $\ge 3$ has eccentricity
$\le 2$}, and let
\[
  \mathrm{RES} := \{\, C : C \text{ connected}, C_4\text{-free},
  \rad(C) = 2,\ \diam(C) = 3,\ \ell(C) > 3,\ \text{(R3b)} \,\}.
\]
Write (C6) for the property that $C$ has no induced $P_6$, i.e.\ $\pth(C) \le 5$.
\end{definition}

Two remarks on the shape of this definition, both of which cost the campaign a round and are
recorded so that the reader does not repeat them.

\begin{remark}[the hypothesis (R3b) is not optional]\label{rem:live}
\MV{w133\_res3\_search.py}
The class $\{C_4$-free, $\rad = 2$, $\diam = 3$, $\ell > 3\}$ --- i.e.\ RES without (R3b) ---
is \emph{inhabited}, with at least nine explicit members: Hoffman--Singleton minus $1,2,3,5,8$
or $12$ vertices ($\ell$ from $6.857$ down to $5.368$) and $ER_5, ER_7$ minus one vertex
($\ell = 3.900, 4.875$). All satisfy the conjecture, with induced-path certificates $\ge 15$.
So no contradiction can be derived from that class as stated. Every one of the nine violates
(R3b), for a structural reason: once a dense core has diameter $3$, the pair realising the
diameter consists of high vertices, which (R3b) forbids.
\end{remark}

\begin{remark}[which reading of (R3)]\label{rem:r3a}
\MV{w133\_res\_kill.py}
The weaker reading (R3a) --- \emph{no pair at distance $\ge 3$ with $a$-values
$(\ge 3, \ge 2)$} --- is \emph{not} equivalent to (R3b), and the class it defines is
inhabited: Hoffman--Singleton with one pendant vertex, with one triangle-leaf, and $ER_q$
with one pendant vertex ($q = 5,7,11$) all satisfy (R3a) and fail (R3b), because the unique
vertex at distance $3$ from the offending high vertex is the grafted $a = 1$ vertex, which
(R3a) never sees. All proofs below use (R3b), which is available only through
Lemma~\ref{lem:g16}.
\end{remark}

\begin{lemma}[core lifting]\label{lem:g16}
\PRVn
Let $G$ be a minimum-order counterexample to \eqref{eq:X} with $\rad(G) = 3$, let
$\Lambda = \{v : a_G(v) = 1\}$ and $C = G - \Lambda$. Then (R3b) holds in $C$.
\end{lemma}

\begin{proof}
By Proposition~\ref{prop:skeleton}(iv), $G$ satisfies $(\star)$ at $r = 3$; more precisely,
$\pth(G) \le \rad(G)+2 = 5$ and Lemma~\ref{lem:g1} give: no pair of $G$ at distance $\ge 3$
has $a_G$-values $(\ge 3, \ge 2)$. Every vertex of $C$ has $a_G \ge 2$ by the definition of
$\Lambda$. Let $h \in C$ with $a_C(h) \ge 3$; deleting vertices cannot increase $a$-values,
so $a_G(h) \ge 3$. If $\ecc_C(h) = 3$, pick $z \in C$ with
$\mathrm{dist}_C(h,z) = 3 = \mathrm{dist}_G(h,z)$ by isometry
(Proposition~\ref{prop:skeleton}(iii)); then $a_G(z) \ge 2$ and $(h,z)$ is a forbidden pair.
\end{proof}

\begin{proposition}[the reduction]\label{prop:reduction}
\PRVn
Let $G$ be a minimum-order counterexample to \eqref{eq:X} with $\rad(G) = 3$. Then the core
$C = G - \Lambda$ belongs to $\mathrm{RES}$ and satisfies (C6).
\end{proposition}

\begin{proof}
By Proposition~\ref{prop:skeleton}, $C$ is connected, $C_4$-free, isometric in $G$, with
$\ell(C) > 3$ and all $a_G$-values $\ge 2$. If $C$ contained a pair at distance $\ge 4$, both
its ends would have $a_C \ge 1$; more usefully, both have $a_G \ge 2$ and hence, by isometry
and Lemma~\ref{lem:g9}, $\pth(G) \ge 4+3 = 7 > 5$. So $\diam(C) \le 3$. If $\diam(C) = 2$,
Theorem~\ref{thm:g12} gives $\pth(G) \ge \pth(C) \ge 6$, a contradiction. So $\diam(C) = 3$.
If $\rad(C) = 3$, minimality applied to $C$ (which has $\ell(C) > 3$ and fewer vertices)
gives $\pth(C) \ge \rad(C)+3 = 6$, again a contradiction; so $\rad(C) = 2$. Property (R3b)
is Lemma~\ref{lem:g16}, and $\pth(C) \le \pth(G) \le 5$ is (C6).
\end{proof}

Theorem~\ref{thm:r3} therefore follows from Theorems~\ref{thm:g26} and \ref{thm:g35} below,
which together say $\mathrm{RES} \wedge (C6) = \emptyset$: with
Corollary~\ref{cor:g8} disposing of $\rad = 2$ and Proposition~\ref{prop:skeleton}(i) of
$\rad \le 1$, a minimum-order counterexample must have radius at least $4$.

Throughout \S\ref{sec:reduction}--\S\ref{sec:hexagon} we fix $C \in \mathrm{RES}$,
$H := \{v : a(v) \ge 3\}$ (all $a$-values are computed in $C$), and we use the following
mass inequality repeatedly.

\begin{lemma}[the mass inequalities]\label{lem:mass}
\PRV{w133\_res\_kill.py}
Let $C$ be $C_4$-free with $\ell(C) > 3$.
\begin{enumerate}
\item[(a)] If every vertex outside $H$ has $a \le 2$ --- which holds by definition of $H$ ---
then $\sum_{h \in H}(a(h)-2) > n$. In particular $H \ne \emptyset$.
\item[(b)] If $c$ has $\ecc(c) = 2$, $A := N(c)$, $k := |A|$ and $B$ is the distance-$2$
layer, then $\sum_{b \in B} a(b) > 3 + k + 2|B|$; consequently $B$ contains a vertex of
$H$.
\end{enumerate}
\end{lemma}

\begin{proof}
(a) $3n < \sum_v a(v) \le \sum_H a(h) + 2(n - |H|)$.
(b) As in the proof of Theorem~\ref{thm:g7}, every $b \in B$ has a unique parent in $A$,
$a(c) = \comp(A) \le k$, and $a(\alpha) = 1 + \comp(B_\alpha)$ for $\alpha \in A$, with
$\sum_\alpha \comp(B_\alpha) \le |B|$. Hence
$\sum_v a(v) \le k + (k + |B|) + \sum_B a(b)$, while $\sum_v a(v) > 3n = 3(1+k+|B|)$. The
consequence follows since $a(b) \le 2$ for all $b \in B$ would give
$\sum_B a(b) \le 2|B|$.
\end{proof}

\begin{lemma}[mass--Bonferroni]\label{lem:g18}
\PRV{w133\_res\_kill.py}
In any $C_4$-free graph and any $S \subseteq V$,
$\sum_{h \in S} a(h) \le |N(S)| + P(S) \le n + \binom{|S|}{2}$, where
$N(S) = \bigcup_{h\in S}N(h)$ and $P(S)$ is the number of pairs of $S$ with a common
neighbour. Consequently, in $\mathrm{RES}$, $|H| \ge 6$.
\end{lemma}

\begin{proof}
Write $a(h) = |I_h|$ for a maximum independent set $I_h \subseteq N(h)$, and
$m(v) := |\{h \in S : v \in I_h\}| \le \deg_S(v)$. Then
$\sum_S a = \sum_v m(v) \le |N(S)| + \sum_v (m(v)-1)^+ \le |N(S)| + \sum_v
\binom{\deg_S(v)}{2}$, and the last sum counts pairs of $S$ with a common neighbour, with
multiplicity equal to the number of common neighbours, which is $\le 1$ by $C_4$-freeness.
For the consequence: Lemma~\ref{lem:mass}(a) gives $\sum_H a > n + 2|H|$, so
$n + \binom{|H|}{2} > n + 2|H|$, i.e.\ $|H| \ge 6$.
\end{proof}

\begin{lemma}[no leaf at eccentricity 3]\label{lem:g19}
\PRV{w133\_res\_kill.py}
In $\mathrm{RES}$ no vertex of eccentricity $3$ is a leaf.
\end{lemma}

\begin{proof}
Let $u$ be a leaf with neighbour $p$ and $\ecc(u) = 3$. Then
$\mathrm{dist}(u,\cdot) = 1 + \mathrm{dist}(p,\cdot) \le 3$ gives $\ecc(p) \le 2$, and
$\ecc(p) = 1$ would force $\diam \le 2$; so $\ecc(p) = 2$. By (R3b) every $h \in H$ has
$\mathrm{dist}(h,u) \le 2$, hence $\mathrm{dist}(h,p) \le 1$, i.e.\ no vertex of $H$ lies at
distance exactly $2$ from $p$ --- contradicting Lemma~\ref{lem:mass}(b) at $c = p$.
\end{proof}

\subsection{The triangular endpoint branch is empty}\label{sec:tri}

By (R3b), a pair $(u,z)$ realising $\diam(C) = 3$ has $a(u), a(z) \le 2$. This subsection
kills the case $a(u) = 1$; \S\ref{sec:hexagon} kills $a(u) = a(z) = 2$.

By Lemma~\ref{lem:g2} and Lemma~\ref{lem:g19}, a vertex $u$ with $a(u) = 1$ and
$\ecc(u) = 3$ has $d(u) = 2$ and $N(u) = \{p,q\}$ with $p \sim q$: a \emph{triangular
endpoint}. Fix such a $u$ and write
\[
  B_p := N(p)\setminus\{u,q\},\quad B_q := N(q)\setminus\{u,p\},\quad
  D := \text{the distance-}3\text{ layer from } u,
\]
$\beta_x := |B_x|$, $\delta_h := |N(h)\cap D|$, $H' := H \cap (B_p \cup B_q)$,
$\eta_x := |H \cap B_x|$, $e_x := $ the number of edges inside $B_x$.

\begin{lemma}\label{lem:g23}
\PRV{w133\_r7\_triend.py}
With this notation:
\begin{enumerate}
\item[(a)] $N(p)\cap N(q) = \{u\}$;
\item[(b)] $\ball_2(u) = \{u,p,q\} \sqcup B_p \sqcup B_q$ and
$|\ball_2(u)| = d(p)+d(q)-1$; in particular $n = 3 + \beta_p + \beta_q + |D|$ and neither $p$
nor $q$ has a neighbour in $D$;
\item[(c)] no vertex of $B_p$ is adjacent to $q$, every vertex has at most one neighbour in
$B_p$ and at most one in $B_q$;
\item[(d)] $a(p) - 2 = \beta_p - 1 - e_p$ and $a(q) - 2 = \beta_q - 1 - e_q$;
\item[(e)] $H \subseteq \{p,q\}\cup B_p \cup B_q$.
\end{enumerate}
\end{lemma}

\begin{proof}
(a) $N(u) = \{p,q\}$ is a clique, and two common neighbours of $p,q$ besides $u$ would give a
$C_4$. (b) Every $b \in \ball_2(u)\setminus\{u,p,q\}$ has exactly one neighbour in $\{p,q\}$,
at least one by definition and at most one by (a); the count is $3 + (d(p)-2)+(d(q)-2)$. A
neighbour of $p$ in $D$ would be at distance $2$ from $u$. (c) $h \in B_p$ adjacent to $q$
would lie in $N(p)\cap N(q) = \{u\}$; two neighbours of a vertex $w$ inside $B_p$ would give
$w$ and $p$ two common neighbours. (d) $N(p) = \{u,q\}\cup B_p$ and its internal edges are
$uq$ together with the $e_p$ edges of $B_p$: there are no $q$--$B_p$ edges by (c) and no
$u$--$B_p$ edges since $N(u) = \{p,q\}$. So $t(p) = 1 + e_p$, $d(p) = 2 + \beta_p$ and
Lemma~\ref{lem:matching} gives $a(p) = 1 + \beta_p - e_p$. (e) By (R3b) every $h \in H$ has
$\ecc(h)\le 2$, so $H \subseteq \ball_2(u)$; and $u \notin H$.
\end{proof}

Note that (d) forces $a(p) \ge 2$ as soon as $\beta_p \ge 1$, because
$e_p \le \lfloor \beta_p/2\rfloor \le \beta_p - 1$ for $\beta_p \ge 2$ and $e_p = 0$ for
$\beta_p = 1$. This is used below to add the $p$- and $q$-terms to the mass sum whether or not
they are hubs.

\begin{lemma}[$P_k$-endpoint degree bound]\label{lem:g22}
\PRV{w133\_r6\_adjudicate.py}
Let $G$ be $C_4$-free with no induced $P_{k+1}$, $k \ge 3$. Then every endpoint of an induced
$P_k$ has degree $\le k$. In particular, $C_4$-free with no induced $P_6$ implies every
endpoint of an induced $P_5$ has degree $\le 5$.
\end{lemma}

\begin{proof}
Let $(v_0,\dots,v_{k-1})$ be an induced $P_k$ and $w \in N(v_0)\setminus\{v_1,\dots,
v_{k-1}\}$. If $w$ had no neighbour among $v_1,\dots,v_{k-1}$ then
$(w,v_0,\dots,v_{k-1})$ would be an induced $P_{k+1}$. So $w \in N(v_0)\cap N(v_i)$ for some
$i \ge 1$, and each such intersection has size $\le 1$ by $C_4$-freeness: that is at most
$k-1$ vertices, and adding $v_1$ gives $d(v_0)\le k$.
\end{proof}

The hypothesis ``endpoint of an induced $P_5$'' cannot be dropped: the friendship graphs
$F_k$ are $C_4$-free and $P_6$-free with unbounded degree $2k$, and contain no induced $P_5$
at all. \MV{w133\_r6\_adjudicate.py}

\begin{lemma}\label{lem:g25}
\REV{w133\_r7\_triend.py}
In $\mathrm{RES}$ with a triangular endpoint $u$:
\begin{enumerate}
\item[(a)] there is no edge between $B_p$ and $B_q$;
\item[(b)] every $h \in B_p \cup B_q$ satisfies $a(h)\le 1 + \delta_h$;
\item[(c)] if $b \in B_p$ has $\ecc(b) = 2$ then $\beta_q \le \delta_b$, and symmetrically;
\item[(d)] $\beta_p, \beta_q \ge 1$;
\item[(e)] under (C6), every hub of $B_p \cup B_q$ has degree $\le 5$, hence $\delta_b \le 4$
for such a hub; and with (c), $\beta_p,\beta_q \le 4$ as soon as both sides carry a hub.
\end{enumerate}
\end{lemma}

\begin{proof}
(a) If $h \in B_p$, $c \in B_q$ and $h \sim c$, then $p$ and $c$ have the two common
neighbours $h$ and $q$ (using $p \sim q$).

(b) $|N(h)\cap B_p| \le 1$ by Lemma~\ref{lem:g23}(c). If $N(h)\cap B_p = \{b\}$ then $pb$ is
an edge inside $N(h)$, so $a(h)\le d(h)-1 = 1 + \delta_h$ using (a); otherwise
$d(h) = 1 + \delta_h \ge a(h)$.

(c) For $c \in B_q$ we have $\mathrm{dist}(b,c) \ge 2$ by (a), so $\ecc(b) = 2$ gives a
common neighbour $w$ of $b$ and $c$. Now $w \ne p$ (as $p \not\sim c$ by
Lemma~\ref{lem:g23}(a)), $w \notin B_p$ by (a), and $w \ne u,q$; so $w \in N(b)\cap D$. Each
$D$-vertex has at most one neighbour in $B_q$, so $c \mapsto w$ is injective.

(d) Suppose $B_q = \emptyset$. Then $N(q) = \{u,p\}$ with $u \sim p$, so $a(q) = 1$ and
$q \notin H$; every $y \in D$ has at most one neighbour in $B_p \cup B_q = B_p$. From
$a(p) \le d(p)-1 = \beta_p + 1$ and $a(h) - 2 \le \delta_h$ for $h \in H'$ (which is (b)
weakened), Lemma~\ref{lem:mass}(a) gives
\[
  n < \sum_{h\in H}(a(h)-2) \le (\beta_p - 1) + \sum_{h\in H'}\delta_h
    \le (\beta_p-1) + |D| = (n-3)-1 = n-4,
\]
a contradiction.

(e) Let $h \in H\cap B_p$ with $d(h)\ge 6$. By (R3b), $\ecc(h) \le 2$, and $\ecc(h) = 1$
would force $\diam \le 2$; so $\ecc(h) = 2$. By (d), $\beta_q \ge 1$. If some
$y \in N(h)\cap D$ had a neighbour $c \in B_q$, then $c \notin N(h)$ by (a) and
$(h,y,c,q,u)$ is an induced $P_5$ --- $h \not\sim c$, $h \not\sim q$, $h \not\sim u$,
$y \not\sim q, u$ (as $y \in D$), $c \not\sim u$ --- so Lemma~\ref{lem:g22} gives
$d(h)\le 5$, a contradiction. Hence no vertex of $N(h)\cap D$ has a $B_q$-neighbour and the
injection of (c) has empty target, giving $\beta_q \le 0$: contradiction. The last clause is
(c) together with $\delta_b \le d(b) - 1 \le 4$.
\end{proof}

\begin{theorem}[the triangular branch is empty]\label{thm:g26}
\REV{w133\_r7\_triend.py}
In $\mathrm{RES} \wedge (C6)$ no vertex has $a$-value $1$ and eccentricity $3$. Equivalently,
every pair of vertices at distance $3$ has both $a$-values equal to $2$.
\end{theorem}

\begin{proof}
Suppose $u$ is such a vertex and keep the notation above. Let
$D_2 := \{y \in D : y$ has a neighbour in $H'\cap B_p$ and one in $H'\cap B_q\}$.

\emph{Step 1: a lower bound on $|D_2|$.} By Lemma~\ref{lem:g23}(d),(e) and
Lemma~\ref{lem:g25}(b), and since $a(p)-2, a(q)-2 \ge 0$ by Lemma~\ref{lem:g25}(d),
\[
 n \;<\; \sum_{h\in H}(a(h)-2)
   \;\le\; (\beta_p - 1 - e_p) + (\beta_q - 1 - e_q) + \sum_{h \in H'} (\delta_h - 1).
\]
Each $y \in D$ has at most one neighbour in $B_p$ and one in $B_q$, so it contributes at most
$1$ to $\sum_{h\in H'}\delta_h$, and $2$ exactly when $y \in D_2$; hence
$\sum_{h\in H'}\delta_h = \sum_{y\in D}|N(y)\cap H'| \le |D| + |D_2|$. Substituting
$n = 3+\beta_p+\beta_q+|D|$ and $|H'| = \eta_p + \eta_q$ yields
\[
  |D_2| \;\ge\; 6 + \eta_p + \eta_q + e_p + e_q .
\]
Two vertices of $D_2$ with the same pair of parents would be two common neighbours of that
pair, a $C_4$; so $y \mapsto (\text{its }B_p\text{-parent}, \text{its }B_q\text{-parent})$ is
injective and $|D_2| \le \eta_p\eta_q$. Therefore
\[
  (\eta_p - 1)(\eta_q - 1) \;\ge\; 7 + e_p + e_q .
\]

\emph{Step 2: both sides are saturated.} The displayed inequality forces
$\eta_p, \eta_q \ge 2$, so both sides carry hubs and Lemma~\ref{lem:g25}(e) gives
$\beta_p,\beta_q \le 4$, hence $\eta_p,\eta_q \le 4$. The only integer pair with both entries
$\le 4$ and $(\eta_p-1)(\eta_q-1)\ge 7$ is $\eta_p = \eta_q = 4$; so $\beta_p = \beta_q = 4$,
\emph{every} vertex of $B_p\cup B_q$ is a hub, and $d(p) = d(q) = 6$.

\emph{Step 3: the $B$-vertices are pinned.} For a hub $b \in B_p$, Lemma~\ref{lem:g25}(c)
gives $\delta_b \ge \beta_q = 4$ while Lemma~\ref{lem:g25}(e) gives
$5 \ge d(b) = 1 + |N(b)\cap B_p| + \delta_b$. Hence $|N(b)\cap B_p| = 0$, $\delta_b = 4$,
$d(b) = 5$; therefore $e_p = e_q = 0$ and $a(p) = a(q) = 5$.

\emph{Step 4: the graph is determined.} $a(p) = 5 \ge 3$ makes $p$ a hub, so $\ecc(p) = 2$
and every $y \in D$ has a $B_p$-parent, and symmetrically a $B_q$-parent. Counting
$B_p$--$D$ edges gives $|D| = \sum_{b\in B_p}\delta_b = 16$, and injectivity makes
$y \mapsto (b,c)$ a bijection $D \to B_p\times B_q$. Thus $n = 3+4+4+16 = 27$.

\emph{Step 5: the contradiction.} Fix $y \in D$ with $B_p$-parent $b_1$ and let $b$ be one of
the other three vertices of $B_p$. As $b$ is a hub, $\ecc(b) = 2$, so $b$ and $y$ have a
common neighbour $w$; $w \ne b_1$ since $e_p = 0$, and $w \notin B_q$ by
Lemma~\ref{lem:g25}(a), so $w \in N(y)\cap D$ --- and the unique $B_p$-parent of $w$ is $b$,
so the three choices of $b$ give three distinct $w$. Hence $d(y) \ge 2+3 = 5$ and, in a
$C_4$-free graph, $a(y)\ge \lceil d(y)/2\rceil \ge 3$: $y$ is a hub at distance $3$ from
$u$, contradicting $H \subseteq \ball_2(u)$.
\end{proof}

\begin{corollary}\label{cor:g26.1}
\REV{w133\_r7\_triend.py}
In $\mathrm{RES}\wedge(C6)$ every pair $(u,z)$ at distance $3$ satisfies $a(u)=a(z)=2$.
\end{corollary}

\begin{remark}
The proof of Theorem~\ref{thm:g26} never uses the branch hypothesis that \emph{every}
diameter pair has a low endpoint --- only the existence of one vertex of $a$-value $1$ at
eccentricity $3$ --- which is why it yields the stronger corollary. Property (C6) enters only
through Lemma~\ref{lem:g25}(e), via Lemma~\ref{lem:g22}. No exhaustive search over graphs is
involved: the only sweeps are the small-graph table backing the unconditional lemmas and two
bounded integer sweeps over $(\eta_p,\eta_q,e_p,e_q) \in [0,4]^4$.
\end{remark}

\subsection{The hexagon branch is empty}\label{sec:hexagon}

By Corollary~\ref{cor:g26.1} the diameter pair $(u,z)$ of $C \in \mathrm{RES}\wedge(C6)$ has
$a(u) = a(z) = 2$. The next lemma turns that into an induced hexagon.

\begin{lemma}[the forced hexagon]\label{lem:g13}
\PRV{wowii133\_gplus.c}
Let $G$ be $C_4$-free with $\mathrm{dist}(u,z) = 3$ and $a(u) = a(z) = 2$. Then either
$\pth(G)\ge 6$, or: $N(u)$ consists of $u_1$'s component together with a \emph{singleton}
component $\{x\}$, $N(z)$ consists of $u_2$'s component together with a singleton $\{y\}$,
and $x \sim y$ --- so that $G$ contains an induced $6$-cycle through $u$ and $z$.
\end{lemma}

\begin{proof}
Fix a geodesic $u, u_1, u_2, z$. Take $x \in N(u)$ outside $u_1$'s component and
$y \in N(z)$ outside $u_2$'s component. By the checks in Lemma~\ref{lem:g10} applied at both
ends, $x,y \notin \{u,u_1,u_2,z\}$ and $x \not\sim u_1,u_2,z$, $y \not\sim u_2,u_1,u$. If
some legitimate choice has $x\not\sim y$, then $x,u,u_1,u_2,z,y$ is an induced $P_6$. So
suppose every legitimate $x$ is adjacent to every legitimate $y$. If $x$'s component had two
vertices $x, \bar x$, then $x$ and $\bar x$ would both be adjacent to $u$ and to $y$, a
$C_4$; so it is a singleton, and symmetrically for $y$.
\end{proof}

So $C$ carries an induced hexagon $Z = (u, u_1, u_2, z, y, x)$ at cyclic positions $0..5$,
with $\{x\}$ a singleton component of $N(u)$ and $\{y\}$ a singleton component of $N(z)$.
The next lemma is the engine of the whole collapse and is of independent interest.

\begin{lemma}\label{lem:g32p}
\REV{w133\_r8\_q61.py}
Let $Z = (v_0,\dots,v_5)$ be an induced $C_6$ in a graph with no induced $P_6$. Then no
vertex outside $Z$ has exactly two $Z$-neighbours that are \emph{consecutive} on $Z$. (No
$C_4$-freeness and no other hypothesis is used.)
\end{lemma}

\begin{proof}
Suppose $w \notin Z$ has exactly the two $Z$-neighbours $v_i, v_{i+1}$. Deleting $v_{i+1}$
from $Z$ leaves the induced path $(v_i, v_{i-1}, v_{i-2}, v_{i-3}, v_{i-4})$ on the remaining
five vertices, with endpoints $v_i$ and $v_{i-4} = v_{i+2}$; and $w$ is adjacent to $v_i$
only among them, its other $Z$-neighbour having been deleted. Hence
$(w, v_i, v_{i-1}, v_{i-2}, v_{i-3}, v_{i-4})$ is an induced $P_6$.
\end{proof}

\begin{lemma}[antipodal rigidity]\label{lem:g32}
\REV{w133\_r8\_q61.py}
Let $Z = (v_0,\dots,v_5)$ be an induced $C_6$ in a $C_4$-free graph with no induced $P_6$.
Then every vertex outside $Z$ has $0$ or $2$ neighbours on $Z$, and if $2$ they are
\emph{antipodal} (cyclic distance exactly $3$).
\end{lemma}

\begin{proof}
Let $w \notin Z$.

\emph{$|N(w)\cap Z| \ne 1$}: if $v_i$ is the unique $Z$-neighbour, delete $v_{i+1}$; the
resulting induced $P_5$ has $v_i$ as an endpoint and $w$ extends it to an induced $P_6$.

\emph{$|N(w)\cap Z| \le 2$}: any three vertices of a hexagon contain a pair at cyclic
distance $2$. Indeed three vertices split the cycle into arcs of lengths $a+b+c = 6$ with
$a,b,c \ge 1$, and a pair separated by an arc of length $j$ is at cyclic distance
$\min(j,6-j)$, which equals $2$ exactly when $j \in \{2,4\}$; if no pair were at cyclic
distance $2$ we would need $a,b,c \in \{1,3,5\}$ with $a+b+c = 6$, which has no solution.
For such a pair $v_j, v_{j+2}$, the vertices $w$ and $v_{j+1}$ have the two common
neighbours $v_j, v_{j+2}$, a $C_4$.

\emph{Cyclic distance $2$ is excluded} by the same $C_4$, and \emph{cyclic distance $1$} by
Lemma~\ref{lem:g32p}.
\end{proof}

\begin{remark}[a statement that outran its proof]\label{rem:g32false}
\MV{w133\_r8\_judge20\_g32.py}
Lemma~\ref{lem:g32} is \emph{false} without $C_4$-freeness: $C_6$ together with a vertex $w$
adjacent to $v_0$ and $v_2$ is $P_6$-free, and $w$ has two $Z$-neighbours at cyclic distance
$2$. The first version of this lemma in the campaign omitted the hypothesis and was refuted by
the blind reviewer; the genuinely hypothesis-free half is Lemma~\ref{lem:g32p}. The defect was
inherited: an upstream lemma about induced hexagons had been \emph{stated} with the
conclusion ``$0$ or $2$ neighbours'' while its written proof only excluded ``exactly $1$'',
never ``$\ge 3$'', and the present Lemma~\ref{lem:g32} copied the half proof together with the
full statement. The missing line --- the arc argument above --- consumes only
$C_4$-freeness. A sweep of all $12$ places where the upstream lemma is used confirmed that
every one carries $C_4$-freeness, so no downstream conclusion moved.
\end{remark}

\begin{corollary}[mass identity]\label{cor:g32.1}
\REV{w133\_r8\_q61.py}
In a $C_4$-free $P_6$-free graph with an induced $C_6$, $t(v_i) = 0$ and
\[
  a(v_i) = d(v_i) = 2 + [\,\text{the } (v_i,v_{i+3}) \text{ slot is occupied}\,].
\]
\end{corollary}

\begin{proof}
By Lemma~\ref{lem:g32} the only off-$Z$ neighbour a cycle vertex $v_i$ can have is the vertex
occupying the antipodal slot $(v_i,v_{i+3})$, of which there is at most one by
$C_4$-freeness. $N(v_i) = \{v_{i-1},v_{i+1}\}$ possibly with that vertex; $v_{i-1}\not\sim
v_{i+1}$ (cyclic distance $2$ on an induced hexagon), and the antipodal-slot vertex is
adjacent to neither. So $N(v_i)$ is independent, $t(v_i)=0$, and Lemma~\ref{lem:matching}
gives $a = d$.
\end{proof}

Over the exhaustive $C_4$-free table on $6 \le n \le 7$ ($111\,508$ connected graphs) plus
$47$ explicit and greedily generated carriers, restricted to the $P_6$-free members carrying
an induced $C_6$, the classification of Lemma~\ref{lem:g32} and the identity of
Corollary~\ref{cor:g32.1} were asserted on $1\,330$ hexagons, $1\,300$ off-hexagon vertices
and $15\,960$ mass identities with $0$ failures; ten hexagons carry two occupied antipodal
slots simultaneously, so the occupied case is tested non-vacuously. \MV{w133\_r8\_q61.py}

\begin{lemma}[the frame collapses]\label{lem:g33}
\REV{w133\_r8\_q61.py}
Let $C \in \mathrm{RES}\wedge(C6)$ with the hexagon $Z = (u,u_1,u_2,z,y,x)$ of
Lemma~\ref{lem:g13}. Then
\[
 N(u) = \{u_1, x\},\quad N(z) = \{u_2,y\},\quad d(u)=d(z)=2,\quad
 V(C) = Z \sqcup \{P?\} \sqcup \{Q?\} \sqcup W,
\]
where $P$ is the (at most one) vertex adjacent to $u_1$ and $y$, $Q$ the (at most one) vertex
adjacent to $u_2$ and $x$, and $W$ the set of vertices with no $Z$-neighbour. Moreover every
$w \in W$ is at distance $3$ from both $u$ and $z$, hence $\ecc(w) = 3$, hence $a(w) = 2$ by
(R3b) and Corollary~\ref{cor:g26.1}: no vertex of $W$ is a hub.
\end{lemma}

\begin{proof}
The antipodal pairs of $Z$ are $(u,z)$, $(u_1,y)$ and $(u_2,x)$. By Lemma~\ref{lem:g32},
every off-$Z$ vertex with a $Z$-neighbour realises one of these three pairs, and $(u,z)$ is
impossible since $\mathrm{dist}(u,z) = 3$; each of the other two is realised by at most one
vertex, since two would be two common neighbours of the pair. Call them $P$ and $Q$. In
particular no off-$Z$ vertex is adjacent to $u$ or to $z$, so $N(u) = N(u)\cap Z = \{u_1,x\}$
and $N(z) = \{u_2,y\}$. If $w$ has no $Z$-neighbour then $w$ is at distance $\ge 2$ from $u$
and shares no neighbour with $u$ (both of $N(u)$'s members lie on $Z$), so
$\mathrm{dist}(u,w) = 3$ because $\diam(C) = 3$; same at $z$. Thus $\ecc(w) = 3$, and (R3b)
gives $a(w)\le 2$ while Corollary~\ref{cor:g26.1} applied to the distance-$3$ pair $(u,w)$
gives $a(w) = 2$.
\end{proof}

The four induced $P_6$'s that kill the four consecutive-pair slots explicitly are
$(u_1', u, x, y, z, u_2)$, $(u_2', z, y, x, u, u_1)$, $(A, u_1, u, x, y, z)$ and
$(B, x, u, u_1, u_2, z)$, where $u_1', u_2', A, B$ denote hypothetical occupants of the slots
$(u,u_1), (u_2,z), (u_1,u_2), (y,x)$. Over $176$ $C_4$-free frames ``hexagon $+$ one
consecutive-pair slot $+$ an arbitrary edge pattern on the other five slots'' the named
witness is an induced $P_6$ every time. \MV{w133\_r8\_q61.py}

\begin{lemma}\label{lem:g34}
\REV{w133\_r8\_q61.py}
In $\mathrm{RES}\wedge(C6)$, write $W_P := N(P)\cap W$. Then
\begin{enumerate}
\item[(a)] $N(w) \subseteq \{P,Q\}\cup W_P$ for every $w \in W_P$;
\item[(b)] every $w \in W_P$ is adjacent to $Q$ (so $W_P \ne \emptyset$ forces $Q$ to exist);
\item[(c)] $|W_P|\le 1$, hence $d(P) \le 4$ and $a(P)\le 3$; symmetrically $a(Q)\le 3$.
\end{enumerate}
\end{lemma}

\begin{proof}
(a) $(w,P,u_1,u,x)$ is an induced $P_5$: $w$ has no $Z$-neighbour; $P \not\sim u,x$ since
$P$'s only $Z$-neighbours are $u_1,y$; $u_1\not\sim x$ on the induced hexagon. By (C6) no
induced $P_6$ extends it, so every $v \in N(w)\setminus\{P\}$ is adjacent to one of
$P,u_1,u,x$. Now $v \notin Z$ (as $w$ has no $Z$-neighbour, $v$ could only be a $Z$-vertex at
distance $1$ from $w$, which is excluded); $v \sim P$ puts $v \in W_P\cup\{Q\}$;
$v \sim u_1$ puts $v \in \{P\}$ (the slot $(u_1,u_2)$ is empty); $v \sim u$ is impossible
since $N(u)\subseteq Z$; and $v \sim x$ puts $v \in \{Q\}$ (the slot $(y,x)$ is empty).

(b) If $w \not\sim Q$ then by (a) $N(w) \subseteq \{P\}\cup W_P$, and
$|N(w)\cap W_P| \le 1$ since two $W_P$-neighbours of $w$ would be two common neighbours of
$w$ and $P$. So either $N(w) = \{P\}$, or $N(w) = \{P,w'\}$ with $w' \in W_P \subseteq N(P)$ and hence
$P \sim w'$; either way $N(w)$ is a clique, so $a(w) = 1$, contradicting $a(w) = 2$ from Lemma~\ref{lem:g33}.

(c) Two vertices of $W_P$ would be two common neighbours of $P$ and $Q$ by (b), a $C_4$.
Then $N(P) = \{u_1,y\}\sqcup(\{Q\}\cup W_P)$ with $u_1 \not\sim y$ (antipodal on an induced
hexagon) and neither adjacent to anything else in $N(P)$, while $Q \sim w$; so
$a(P) \le 2 + 1 = 3$.
\end{proof}

\begin{theorem}[the hexagon branch is empty]\label{thm:g35}
\REV{w133\_r8\_q61.py, w133\_r9\_frame\_l.py}
$\mathrm{RES}\wedge(C6) = \emptyset$. Equivalently: every connected $C_4$-free graph with
$\rad = 2$, $\diam = 3$, $\ell > 3$ and (R3b) contains an induced $P_6$.
\end{theorem}

\begin{proof}
Suppose $C \in \mathrm{RES}\wedge (C6)$ and set up the frame above. By
Lemma~\ref{lem:g33} every hub of $C$ lies in $\{u_1,u_2,x,y,P,Q\}$: the vertices $u,z$ have
$a = 2$, no vertex of $W$ is a hub, and there is nothing else. By
Corollary~\ref{cor:g32.1} and Lemma~\ref{lem:g34},
\[
  a(u_1)-2 = [P],\quad a(u_2)-2 = [Q],\quad a(x)-2=[Q],\quad a(y)-2=[P],
\]
\[
  a(P)-2 \le 1,\qquad a(Q)-2\le 1,
\]
where $[P],[Q] \in \{0,1\}$ record whether the slots are occupied. Hence
\[
  \sum_{h\in H}\bigl(a(h)-2\bigr) \;\le\; 3[P]+3[Q] \;\le\; 6 .
\]
But Lemma~\ref{lem:mass}(a) demands $\sum_{h\in H}(a(h)-2) > n$, and $n \ge |Z| = 6$, so the
mass must be at least $7$: a contradiction.
\end{proof}

\begin{remark}[the margin, and a second kill]\label{rem:secondkill}
\REV{w133\_r9\_frame\_l.py}
The margin is not thin. By Lemma~\ref{lem:g18}, $|H|\ge 6$, which forces all six candidate
slots to be hubs, hence $[P]=[Q]=1$, hence $n \ge 8$ and a required mass of $\ge 9$ against a
cap of exactly $6$.

Moreover the class dies a second time, by a route that never uses the mass inequality. By
Lemma~\ref{lem:g33}, $V = Z\sqcup\{P?\}\sqcup\{Q?\}\sqcup W$ and $N(u) = \{u_1,x\}\subseteq
Z$. For $w \in W$, $\mathrm{dist}(w,u) = 3$, so a shortest $w$--$u$ path is $w,a,b,u$ with
$b \in \{u_1,x\}$ and $a \in N(u_1)\cup N(x) = \{u,u_2,P\}\cup\{u,y,Q\}$; adjacency to $w$
forces $a \notin Z$, i.e.\ $a \in \{P,Q\}$. So $W = W_P\cup W_Q$, and
Lemma~\ref{lem:g34}(b) gives $W_P \subseteq W_Q$ and $W_Q\subseteq W_P$, whence $W = W_P$ and
$|W| \le 1$ and $n = 6 + [P]+[Q]+|W| \le 9$. The enumeration of all frames
$Z \sqcup\{P?\}\sqcup\{Q?\}\sqcup W$ with $|W|\le 3$ and all edge patterns is therefore
\emph{exhaustive} for the class; it returns exactly $5$ admissible frames, with
\[
 (P,Q,|W|{=}1):\; n=9,\ \ell = 8/3; \qquad (P,Q):\; n=8,\ \ell=2.5;
\]
\[
 (P) \text{ or } (Q):\; n=7,\ \ell=16/7; \qquad (\text{neither}):\; n=6,\ \ell=2 .
\]
So $\max \ell = 8/3 < 3$, contradicting $\ell > 3$ directly. Both kills depend on
\S\ref{sec:tri}; exactly one of them depends on the mass inequality.
\end{remark}

\begin{corollary}\label{cor:g35.1}
\REV{w133\_r8\_q61.py}
$C \in \mathrm{RES}$ implies $\pth(C)\ge 6$.
\end{corollary}

\begin{proof}
Theorems~\ref{thm:g26} and \ref{thm:g35}.
\end{proof}

\begin{proof}[Proof of Theorem~\ref{thm:r3}]
Let $G$ be a minimum-order counterexample to \eqref{eq:X}. By
Proposition~\ref{prop:skeleton}(i), $\rad(G)\ge 2$; $\rad(G) = 2$ is excluded by
Corollary~\ref{cor:g8}. If $\rad(G) = 3$, Proposition~\ref{prop:reduction} puts the core $C$
in $\mathrm{RES}$ with (C6), contradicting Corollary~\ref{cor:g35.1} since
$\pth(G)\ge\pth(C)\ge 6 = \rad(G)+3$.
\end{proof}

\begin{remark}[what is not reproduced]
An earlier route to Theorem~\ref{thm:g35} ran through an explicit list of eleven possible
slots for a hub and reduced the branch to a finite window $14 \le n \le 20$. Lemma~\ref{lem:g32}
supersedes it: the window is closed by proof rather than by search, the intermediate lemmas
become vacuous (their subjects $u_1',u_2',A,B$ do not exist), and the window's lower endpoint
was in any case withdrawn as unsupported by its own logs --- the correct and sufficient bound
is $n \ge |Z| = 6$. None of that material is reproduced here. Likewise a ball-confinement
bound $|H| \le 24$ and a centre identity, both unconditional and both used only by the
superseded route, are omitted.
\end{remark}

\section{The band $\lfloor \ell\rfloor \ge 4$}\label{sec:adopted}

The three results of this section carry the lower \AD{} tier of \S\ref{sec:conventions}: the
arguments originated in the free exploration channel, were re-derived line by line here
(with the two repairs marked below), and passed \emph{one} adversarial review pass, which
returned no defect at any joint classified as mathematics. They are not on the same footing
as \S\ref{sec:r3}.

The first two extend Lemma~\ref{lem:g14} downwards. Recall that its proof used $d \ge 5$
exactly twice, to keep the two extension vertices $w$ and $y$ far apart.

\begin{lemma}[$d = 4$]\label{lem:g36}
\AD{w133\_r10\_adjudicate.py, w133\_r11\_adjudicate.py}
Let $G$ be $C_4$-free with a geodesic $u_0u_1u_2u_3u_4$, $a(u_0)\ge 3$, $a(u_4)\ge 2$, and a
usable side-neighbour $x$ of $u_0$ with $a(x)\ge 4$. Then $\pth(G)\ge 8 = d+4$. No further
hypothesis is needed.
\end{lemma}

\begin{proof}
Choose $y \in N(u_4)$ outside $u_3$'s component, so $y \not\sim u_3$.

\emph{Step 1: the kill accounting.} $N(x)$ has at least $a(x)\ge 4$ components; $u_0$'s is
one, and we let $C_1,C_2,C_3$ be three others with representatives $w_1,w_2,w_3$. No
neighbour of $x$ is adjacent to $u_4$, since
$\mathrm{dist}(w,u_4)\ge \mathrm{dist}(u_0,u_4)-\mathrm{dist}(u_0,w) = 4-2 = 2$ for
$w \in N(x)$. Each of $y, u_2, u_3$ destroys at most one component, because $x$ is adjacent
to none of them --- $x \not\sim u_1$ by usability; $x\not\sim u_2$, else $u_0,u_2$ have the
common neighbours $x,u_1$; $x\not\sim u_3$ by distance;
$\mathrm{dist}(x,y)\ge\mathrm{dist}(u_0,y)-1 \ge 2$ using
$\mathrm{dist}(u_0,y)\ge \mathrm{dist}(u_0,u_4)-1 = 3$ --- and $C_4$-freeness gives
$|N(x)\cap N(v)|\le 1$ for every $v \not\sim x$.

\emph{Step 2: the standard path.} If some $w_i$ avoids all three kills, then
$w_i, x, u_0, u_1, u_2, u_3, u_4, y$ is an induced $P_8$. The chords: $w_i \not\sim u_0$
(different components of $N(x)$); $w_i \not\sim u_1$, else $w_i x u_0 u_1$ is a $C_4$ (using
$x \not\sim u_1$); $w_i \not\sim u_2,u_3$ by assumption and $w_i\not\sim u_4$ by Step~1;
$x \not\sim u_2,u_3,u_4,y$; the geodesic supplies $u_0\not\sim u_2,u_3,u_4$,
$u_1\not\sim u_3,u_4$, $u_2 \not\sim u_4$; $y \not\sim u_3$ by choice; $y \not\sim u_2$, else
$u_2,u_4$ have the common neighbours $y,u_3$; and $y \not\sim u_1,u_0$ by distance. All eight
vertices are distinct: $w_i, x \in \ball_2(u_0)$ with $w_i \ne u_2$ (as $x \not\sim u_2$),
and $y \notin\{u_2,u_3\}$ since $y \not\sim u_3$ while $y \sim u_4$.

\emph{Step 3: all three kills land, on three distinct components.} After relabelling,
$w_y \sim y$, $w_2 \sim u_2$, $w_3 \sim u_3$ lie in three distinct components of $N(x)$, none
of them $u_0$'s. Since $a(u_0)\ge 3$ we may pick $z \in N(u_0)$ outside \emph{both} $u_1$'s
and $x$'s components, so $z \not\sim u_1$ and $z\not\sim x$. Then $z \not\sim y$: otherwise
$\mathrm{dist}(u_0,y)\le 1 + \mathrm{dist}(z,y) = 2$, contradicting
$\mathrm{dist}(u_0,y)\ge 3$. Take
\[
  z - u_0 - x - w_y - y - u_4 - u_3 - u_2 \qquad (8 \text{ vertices}).
\]
Its edges are present by construction. The twenty-one non-adjacent pairs: $z \not\sim x$
(components of $N(u_0)$); $z\not\sim w_y$, else $u_0, w_y$ have the common neighbours $x,z$;
$z \not\sim y$ (above); $z \not\sim u_4, u_3$ by distance; $z \not\sim u_2$, else $u_0,u_2$
have the common neighbours $z,u_1$; $u_0 \not\sim w_y$ (components of $N(x)$);
$u_0 \not\sim y$ by distance; $u_0\not\sim u_4,u_3,u_2$ (geodesic); $x \not\sim y$,
$x\not\sim u_4$ by distance, $x\not\sim u_3,u_2$ (Step 1); $w_y \not\sim u_4$ (Step 1);
$w_y \not\sim u_3$ and $w_y \not\sim u_2$, because a vertex of $N(x)$ adjacent to $u_3$
(resp.\ $u_2$) must lie in the single component that $u_3$ (resp.\ $u_2$) kills, which is
that of $w_3$ (resp.\ $w_2$), and $w_y$ lies in a third; $y \not\sim u_3$ by choice;
$y \not\sim u_2$ as in Step 2; $u_4 \not\sim u_2$ (geodesic). Distinctness: $z \ne x$ and
$z \notin N(x)$; $w_y \in N(x)$ lies at distance $2$ from $u_0$ and $w_y \ne u_2$ since
$x \not\sim u_2$; $y \notin \{u_2,u_3\}$ as in Step 2.
\end{proof}

\begin{lemma}[$d = 3$, Case 1]\label{lem:g37}
\AD{w133\_r10\_adjudicate.py, w133\_r11\_adjudicate.py}
Let $G$ be $C_4$-free with a geodesic $u_0u_1u_2u_3$, $a(u_0)\ge 3$, $a(u_3)\ge 2$, a usable
side-neighbour $x$ of $u_0$ with $a(x)\ge 4$, \emph{and} a usable side-neighbour $y$ of $u_3$
with $y \not\sim x$. Then $\pth(G)\ge 7 = d+4$.
\end{lemma}

\begin{proof}
Fix such a $y$. $N(x)$ has $\ge 4$ components, one of them $u_0$'s; pick representatives
$w_1,w_2,w_3$ of three others. As in Lemma~\ref{lem:g36}, each of $y,u_2,u_3$ kills at most
one component, using $x\not\sim u_2$ (the $C_4$ through $u_0,u_1$), $x\not\sim u_3$
(distance) and $x \not\sim y$ (hypothesis).

\emph{Subcase 1a: some $w_i$ avoids all three kills.} Then
$w_i, x, u_0, u_1, u_2, u_3, y$ is induced, by the chord list of Lemma~\ref{lem:g36} Step 2
read at $d = 3$: $w_i\not\sim u_0$ (component), $w_i \not\sim u_1$ (the $C_4$ through
$x, u_0$), $w_i\not\sim u_2,u_3,y$ (kills avoided), $x \not\sim u_1,u_2,u_3,y$,
$u_0\not\sim u_2,u_3$, $u_1\not\sim u_3$, $y \not\sim u_2$ (choice of $y$), $y\not\sim u_1$
(else $u_1,u_3$ have the common neighbours $y,u_2$), $y \not\sim u_0$ (distance).

\emph{Subcase 1b: the three kills land on three distinct components}, say $w_1 \sim y$,
$w_2\sim u_2$, $w_3\sim u_3$. Since $a(u_0)\ge 3$, pick $z \in N(u_0)$ outside both $u_1$'s
and $x$'s components.

\emph{1b-i: $z \not\sim y$.} Then $z, u_0, x, w_1, y, u_3, u_2$ is induced:
$z\not\sim x$ (components of $N(u_0)$); $z \not\sim w_1$, else $u_0,w_1$ have the common
neighbours $x,z$; $z\not\sim y$ (assumed); $z \not\sim u_3$ (distance); $z \not\sim u_2$,
else $u_0,u_2$ have the common neighbours $z,u_1$; $u_0 \not\sim w_1$ (components of $N(x)$);
$u_0\not\sim y$ (distance); $u_0 \not\sim u_3,u_2$ (geodesic); $x \not\sim y$ (hypothesis);
$x \not\sim u_3$ (distance); $x\not\sim u_2$ ($C_4$); $w_1\not\sim u_2,u_3$ (the $u_2$- and
$u_3$-kills lie in other components); $y \not\sim u_2$ (choice of $y$).

\emph{1b-ii: $z \sim y$.} Then $w_1, y, z, u_0, u_1, u_2, w_2$ is induced. Its edges are
$w_1y$ (the $y$-kill), $yz$ (assumed), $zu_0$, $u_0u_1$, $u_1u_2$, $u_2w_2$ (the $u_2$-kill).
The chords: $w_1\not\sim z$, else $u_0$ and $w_1$ have the two common neighbours $x,z$;
$w_1\not\sim u_0$ (components of $N(x)$); $w_1\not\sim u_1$, else $x,u_1$ have the common
neighbours $w_1,u_0$; $w_1\not\sim u_2$ and $w_1 \not\sim w_2$ (different components of
$N(x)$); $y \not\sim u_0$ (distance); $y\not\sim u_1$ (else $u_1,u_3$ have the common
neighbours $y,u_2$); $y \not\sim u_2$ (choice of $y$); $y \not\sim w_2$ ---
\emph{[first repaired step]} since $x \not\sim y$, $C_4$-freeness gives
$|N(x)\cap N(y)|\le 1$ and $w_1$ already lies in it, so $w_2 \notin N(y)$;
$z\not\sim u_1$ (components of $N(u_0)$); $z \not\sim u_2$ (else $u_0,u_2$ have the common
neighbours $z,u_1$); $z\not\sim w_2$ (else $u_0,w_2$ have the common neighbours $x,z$);
$u_0\not\sim u_2$ (geodesic); $u_0 \not\sim w_2$, since $w_2 \in N(x)$ lies outside $u_0$'s
component of $N(x)$; and $u_1 \not\sim w_2$ --- \emph{[second repaired step]}: the source text
justified this by ``$u_1$ and $w_2$ share $u_2$ and are not adjacent'', which is not a
reason; the correct reason is that $u_1 \sim w_2$ would give $x$ and $u_1$ the two common
neighbours $u_0$ and $w_2$, a $C_4$ (note $x \ne u_1$ and $u_0 \ne w_2$).
\end{proof}

The complementary case at $d = 3$ --- every usable side-neighbour $y$ of $u_3$ satisfies
$y \sim x$ --- is \emph{not} covered by Lemma~\ref{lem:g37}, and \S\ref{sec:limits} shows
that it is false as posed, with and without the natural repair hypothesis.

The third result is global.

\begin{definition}
Let $G$ be connected and $C_4$-free. A geodesic is \emph{$3$-capped at the end $u_0$} if
every usable side-neighbour $x$ of $u_0$ has $a(x)\le 3$ (vacuously so if there is none).
$G$ is \emph{peripherally $3$-capped} if every geodesic of length $\ge \rad(G)$ is $3$-capped
at both of its ends.
\end{definition}

\begin{theorem}\label{thm:g38}
\AD{w133\_r10\_adjudicate.py, w133\_r11\_adjudicate.py}
Every connected $C_4$-free peripherally $3$-capped graph satisfies $\ell(G) < 4$.
\end{theorem}

\begin{proof}
Write $S := \{v : a(v)\ge 4\}$, $T := \{v : a(v)\le 3\}$, $R := \rad(G)$, and suppose for
contradiction $\sum_v a(v) \ge 4n$.

\emph{Claim 1: every vertex $v$ has at most $2$ neighbours in $S$.} Since $\ecc(v)\ge R$
there is $w$ with $\mathrm{dist}(v,w) = \ecc(v)\ge R$; any geodesic $P$ from $v$ to $w$ has
length $\ge R$ and is therefore $3$-capped at $v$. The usable side-neighbours of $v$ with
respect to $P$ are exactly $N(v)\setminus C_0$, where $C_0$ is the component of $P$'s first
step in $N(v)$, and $|C_0|\le 2$ by Lemma~\ref{lem:matching}. An $S$-neighbour has $a \ge 4$
and so cannot be a usable side-neighbour; hence every $S$-neighbour of $v$ lies in $C_0$.

\emph{Claim 2: if $v$ has two $S$-neighbours $s_1,s_2$, then $s_1\sim s_2$.} By Claim~1's
proof both lie in $C_0$, and $|C_0|\le 2$, so $C_0 = \{s_1,s_2\}$ is a matching edge of
$N(v)$.

\emph{Claim 3: $G[S]$ is a disjoint union of triangles, edges and isolated vertices.} By
Claim~1, $\Delta(G[S])\le 2$, so the components are paths and cycles; Claim~2 forces the two
$S$-neighbours of any degree-$2$ vertex of $G[S]$ to be adjacent, and the resulting chord
raises a degree to $3$ unless the component is a triangle or smaller.

Let $k_3,k_2,k_1$ count the $K_3$-, $K_2$- and $K_1$-components of $G[S]$, so
$|S| = 3k_3+2k_2+k_1$.

\emph{Claim 4.} For a $K_3$ component $\{a,b,c\}$, no $T$-vertex is adjacent to two of
$a,b,c$: if $u$ is adjacent to $a$ and $b$, then $a,b$ are two common neighbours of $u$ and
$c$. \emph{Claim 5.} For a $K_2$ component $\{a,b\}$, at most one $T$-vertex is adjacent to
both, for the same reason. \emph{Claim 6.} If $u \in T$ has two $S$-neighbours they are the
endpoints of a $K_2$ component: by Claim~2 applied to $u$ they are adjacent, hence form an
edge of $G[S]$, which lies in a $K_2$ or a $K_3$ component, and $K_3$ is excluded by
Claim~4.

\emph{Upper bound.} Partition $T$ by $d_S(u)\in\{0,1,2\}$ into $T_0,T_1,T_2$. By Claims~5
and 6, $|T_2|\le k_2$, so
$E(S,T) = 2|T_2|+|T_1| \le k_2 + |T| = k_2 + n - |S|$.

\emph{Lower bound.} For $v \in S$ split $t(v)$ into $t_{SS}+t_{ST}+t_{TT}$ by the location of
the pair. If $v$ lies in a $K_3$ then $d_S(v) = 2$ and those two neighbours are matched to
each other in $N(v)$, so $t_{SS}(v) = 1$, $t_{ST}(v) = 0$ and
$a(v) = (2+d_T(v)) - (1+t_{TT}(v))$, giving $d_T(v)\ge a(v)-1$. If $v$ lies in a $K_2$ then
$d_S(v) = 1$, $t_{SS}(v) = 0$ and $t_{ST}(v)\le 1$, giving again $d_T(v)\ge a(v)-1$. If $v$
is a $K_1$ then $d_S(v) = 0$ and $d_T(v)\ge a(v)$. Summing, with $|S| - k_1$ the number of
$S$-vertices in $K_3$- and $K_2$-components,
$E(S,T) = \sum_{v\in S} d_T(v) \ge \sum_{v\in S}a(v) - (|S|-k_1)$. Bounding each $T$-vertex's
$a$-value by $3$ in $\sum_v a(v)\ge 4n$ gives $\sum_{v\in S}a(v)\ge 4n - 3(n-|S|) = n+3|S|$,
hence $E(S,T)\ge n+2|S|+k_1$.

\emph{The contradiction.} Combining, $n+2|S|+k_1 \le k_2+n-|S|$, i.e.\
$3|S| + k_1 \le k_2$, i.e.\ $9k_3+5k_2+4k_1\le 0$, so $k_3=k_2=k_1=0$ and $S = \emptyset$.
But then every $a(v)\le 3$ and $\sum_v a(v)\le 3n < 4n$.
\end{proof}

\begin{corollary}\label{cor:g38}
\AD{w133\_r10\_adjudicate.py}
Every connected $C_4$-free graph with $\ell(G) > 4$ has, at \emph{every} radius, a geodesic
of length $\ge \rad(G)$ with an end $u_0$ carrying a usable side-neighbour $x$ with
$a(x)\ge 4$.
\end{corollary}

The lemma chain of Theorem~\ref{thm:g38} was machine-tested non-vacuously: $449$ peripherally
$3$-capped graphs in the sweep, $60$ of them with $S \ne \emptyset$, with $K_1$-, $K_2$-,
$K_3$-components and a $T_2$ straddler all realised, all passing every claim and both bounds.
Consistency was checked on the two natural near-misses: the incidence graph of $PG(2,3)$
($n = 26$, $4$-regular, $\ell = 4.0$) and that of $PG(3,2)$ ($n = 50$, $\ell = 4.2$) are both
verified \emph{not} peripherally $3$-capped. \MV{w133\_r10\_adjudicate.py}

\begin{remark}[the boundary, and a vacuity annotation]
Run with strict inequalities the chain already yields $9k_3+5k_2+4k_1 < 0$ without the
$S = \emptyset$ endgame; the non-strict run plus that endgame is exactly what buys the
boundary case $\ell = 4$, so the conclusion is $\ell < 4$ and not merely $\ell \le 4$. The
review pass confirmed the boundary run independently. At the time Theorem~\ref{thm:g38} was
recorded, no connected $C_4$-free graph with $\ell > 4$ and $\rad \ge 5$ was known to the
campaign, so the $\rad \ge 5$ slice of Corollary~\ref{cor:g38} carried an explicit
\emph{possibly vacuous} annotation; the theorem itself is not vacuous, since it executes on
the $\rad = 3$ class. That annotation is removed by Proposition~\ref{prop:w2}, which exhibits
such a graph, and the $\rad \ge 5$ slice is taken up in \S\ref{sec:pocket2}.
\end{remark}

\section{The band $\lfloor \ell\rfloor \ge 4$ at large radius}\label{sec:pocket2}

This section treats the slice $\rad(G)\ge 5$ of the band opened in \S\ref{sec:adopted}. Its
numbered contents carry two different tiers, and the difference matters.
Lemma~\ref{lem:g41}, Theorem~\ref{thm:g42}, Lemma~\ref{lem:g43} and
Corollary~\ref{cor:g44} carry the two-family reviewed tier \REV{} of \S\ref{sec:conventions}:
each passed two adversarial review passes by reviewers drawn from two disjoint model
families, both returning no surviving defect at any joint classified as mathematics. That is
one tier above the \AD{} results of the previous section, and it is a statement about
\emph{review}, not about reach. The band $\lfloor\ell\rfloor \ge 4$ remains open, no part of
the conjecture is closed by anything below, nothing here is formalised, and the two witness
statements Proposition~\ref{prop:w2} and Theorem~\ref{thm:g48} are computational and carry
\MV{} only.

\subsection{A hard positive control at radius $5$}

Every bound in \S\ref{sec:adopted} is proved at a peripheral geodesic, so it is worth knowing
that the regime $\rad \ge 5$ is inhabited at all, and inhabited in the range where the cheap
devices fail.

\begin{proposition}[the control $W_2$]\label{prop:w2}
\MV{w133\_r11\_qradp.py, w133\_r15\_key.py, w133\_r17\_w2\_recompute.py}
There is a connected $C_4$-free graph $G$ with $\ell(G) = 6 > 4$, $\rad(G) = \diam(G) = 5$,
and $\pth(G)\ge 168$.
\end{proposition}

The construction is worth isolating, because it is reusable. A vertex-transitive graph has
$\rad = \diam$; a triangle-free graph has $a(v) = d(v)$ for every $v$, so a $k$-regular
triangle-free graph has $\ell = k$ exactly. Hence \emph{any} vertex-transitive $C_4$-free
triangle-free $k$-regular graph with $k \ge 5$ and $\diam \ge 5$ witnesses
Proposition~\ref{prop:w2}, and Cayley graphs of $\mathrm{PSL}(2,p)$ supply them. The standing
control is
\[
  W_2 \;=\; \mathrm{Cay}\bigl(\mathrm{PSL}(2,11),\, S\bigr), \qquad
  S = \{\,g,\,g^{-1} \;:\; g \in \{M_1,M_2,M_3\}\,\},
\]
\[
  M_1 = \begin{pmatrix} 5 & 0\\ 2 & 9\end{pmatrix},\quad
  M_2 = \begin{pmatrix} 1 & 5\\ 0 & 1\end{pmatrix},\quad
  M_3 = \begin{pmatrix} 1 & 6\\ 10 & 6\end{pmatrix}
  \qquad \text{over } \mathrm{GF}(11),
\]
the matrices being taken in $\mathrm{SL}(2,11)$ modulo $\pm I$. Verified: $n = 660$,
$1980$ edges, connected, $6$-regular, triangle-free, $C_4$-free in the strong sense (no two
vertices with two common neighbours), hence $a(v) = 6$ at every vertex and $\ell = 6$;
$\rad = \diam = 5$; and a greedy induced path on $168$ vertices, every one of whose
non-consecutive pairs was certified non-adjacent.

These invariants are \emph{double-sourced}, which not everything in this paper is. They were
computed once from the matrix model above, and once again by an implementation sharing no
code with it, in which $\mathrm{PSL}(2,11)$ is realised instead as a permutation group on the
twelve points of the projective line over $\mathrm{GF}(11)$ under the M\"obius action, its
faithfulness checked rather than assumed ($|{\rm image}| = 660$, not $1320$). The second
implementation reproduces every number above and strengthens two of them: eccentricity $5$ is
verified at \emph{all} $660$ vertices rather than at three sampled ones, and the certified
induced path it finds has $192$ vertices, so $\pth(W_2)\ge 192$. We keep the recorded bound at
$168$, since that is the number the results below were checked against.

$W_2$ is a hard control rather than a decorative one. Here $\diam + 1 = 6$ lies strictly below
$\rad + 4 = 9$, so the trivial diameter device does not settle it, and the conjecture holds on
it with a factor of about $19$ to spare. A dead end follows: no counterexample in this band
can be a vertex-transitive $C_4$-free graph of degree $\ge 5$ with an easily found long
induced path. If tightness exists in this band it lives in graphs of low $a$-value, rich in
triangles and irregular --- which is also where the residual of \S\ref{sec:r3} lives. Larger
members of the same family were verified the same way and behave identically:
$\mathrm{PSL}(2,13)$ ($n = 1092$, $\rad = \diam = 6$, $\pth \ge 239$), $\mathrm{PSL}(2,17)$
and $\mathrm{PSL}(2,19)$ ($\rad = \diam = 7$), and $\mathrm{PSL}(2,23)$ ($n = 6072$,
$\rad = \diam = 7$).

\subsection{The seed hypothesis is redundant at $d \ge 5$}

\begin{lemma}[$d\ge 5$: the seed needs only $a(u_0)\ge 2$]\label{lem:g41}
\REV{w133\_r11\_g14sharp.py}
Let $G$ be $C_4$-free with a geodesic $u_0u_1\cdots u_d$, $d \ge 5$, $a(u_0)\ge 2$,
$a(u_d)\ge 2$, and a usable side-neighbour $x$ of $u_0$ with $a(x)\ge 4$. Then
$\pth(G)\ge d+4$.
\end{lemma}

\begin{proof}
The construction is that of Lemma~\ref{lem:g14}, unchanged. Take $y \in N(u_d)$ outside
$u_{d-1}$'s component of $G[N(u_d)]$, and $w \in N(x)$ in a component of $G[N(x)]$ other than
$u_0$'s and surviving the $u_2$- and $u_3$-kills, for which $a(x)\ge 4$ pays. Then
$w, x, u_0, \dots, u_d, y$ is induced, by exactly the non-adjacencies listed in the proof of
Lemma~\ref{lem:g14}. That construction never selects a \emph{third} component of $G[N(u_0)]$:
it names only $x$. So all it consumes at $u_0$ is $a(u_0)\ge 2$, which the phrase ``some
usable side-neighbour $x$'' already presupposes.
\end{proof}

The redundancy is specific to $d \ge 5$ and does not propagate downwards. At $d = 4$
(Lemma~\ref{lem:g36}) and $d = 3$ (Lemma~\ref{lem:g37}) the proofs choose a vertex of
$N(u_0)$ outside \emph{both} $u_1$'s and $x$'s components, which genuinely needs
$a(u_0)\ge 3$. Lemma~\ref{lem:g41} was machine-checked where it is not vacuous, namely where
the deleted hypothesis actually fails: on chains built from the incidence graph of
$\mathrm{PG}(2,5)$, $120$ frames with $a(u_0) = 2$ exactly and $d \ge 5$, the construction was
built and certified chord-free every time.

\begin{theorem}[$\mu \ge 2$ at radius $\ge 5$]\label{thm:g42}
\REV{w133\_r11\_g42.py, w133\_r15\_key.py}
Let $G$ be connected and $C_4$-free with $\ell(G) > 4$, $\rad(G)\ge 5$ and $\mu(G)\ge 2$.
Then $\pth(G) \ge \rad(G) + 4$.
\end{theorem}

\begin{proof}
By Corollary~\ref{cor:g38}, $G$ is not peripherally $3$-capped: there is a geodesic
$u_0u_1\cdots u_d$ with $d \ge \rad(G)\ge 5$ and a usable side-neighbour $x$ of $u_0$ with
$a(x)\ge 4$. Since $\mu(G)\ge 2$, both $a(u_0)\ge 2$ and $a(u_d)\ge 2$ hold for free, and by
Lemma~\ref{lem:g41} the first of these is all the seed needs at $d \ge 5$. Lemma~\ref{lem:g41}
gives $\pth(G)\ge d + 4 \ge \rad(G)+4$.
\end{proof}

The mechanism was run end to end fifty times on two witnesses --- $W_2$ and a second graph
built from two disjoint copies of the incidence graph of $\mathrm{PG}(2,5)$ joined by an
internal path --- locating a genuine cap-break geodesic on the graph, building the path by the
proof's own recipe, certifying it induced, and comparing it with $\rad + 4$.

\begin{remark}[two numbers, not one]\label{rem:twonumbers}
On $W_2$ the path that this mechanism \emph{constructs} has exactly $9$ vertices, and
$9 = \rad(W_2)+4$; so the construction delivers the bound with no slack at all. The graph's
own longest induced path is a different quantity: $\pth(W_2)\ge 168$. These are two numbers
and they must not be merged. What is tight on $W_2$ is the \emph{mechanism}, and the $9$
measures the mechanism's slack; the graph itself is nowhere near tight, and the $168$ measures
the graph's. No statement of the form ``the bound is tight on $W_2$'' is being made, and none
would be true.
\end{remark}

\subsection{The $a = 1$ case, and what is left of it}

Theorem~\ref{thm:g42} leaves exactly the graphs carrying a vertex of $a$-value $1$. Peeling
those vertices reaches $\mu \ge 2$ and preserves connectivity and $C_4$-freeness
(Lemma~\ref{lem:g6}); it preserves $\ell > 4$ as well, since by Lemma~\ref{lem:g11}'s
bookkeeping deleting a leaf lowers $\Sigma := \sum_u a(u)$ by $2$ and $n$ by $1$ while
deleting a triangle-leaf lowers $\Sigma$ by $1$ and $n$ by $1$, so $\Sigma - 4n$ rises either
way. But peeling may lower the radius, and the peeled graph's $\pth$ is only a lower bound
for the original's. The next two statements attack the far end directly instead.

\begin{lemma}[far-end truncation]\label{lem:g43}
\REV{w133\_r11\_g43.py}
Let $G$ be $C_4$-free with a geodesic $u_0u_1\cdots u_d$, $d \ge 6$, $a(u_0)\ge 2$, and a
usable side-neighbour $x$ of $u_0$ with $a(x)\ge 4$. Then $\pth(G)\ge d+3$, with no hypothesis
whatever on $a(u_d)$.
\end{lemma}

\begin{proof}
If $a(u_d)\ge 2$ then Lemma~\ref{lem:g41} already gives $d+4$. So suppose $a(u_d) = 1$; then
by Lemma~\ref{lem:g2}, $N(u_d)$ is either $\{u_{d-1}\}$ or $\{u_{d-1}, q\}$ with
$q \sim u_{d-1}$. Truncate the geodesic to $u_0u_1\cdots u_{d-1}$, of length $d-1\ge 5$, and
take $y := u_d$ as its far-side vertex. This is legitimate: $u_d \in N(u_{d-1})$, and $u_d$
lies outside $u_{d-2}$'s component of $G[N(u_{d-1})]$. In the first case $u_d$ is isolated in
that graph. In the second its component is $\{u_d,q\}$, and $u_{d-2} \notin \{u_d,q\}$: it is
not $u_d$, which is at distance $2$ from it, and if it were $q$ then $u_{d-2}$ and $u_d$ would
be adjacent, contradicting the geodesic. Nor is $u_{d-2}\sim q$, since $q$ would then have two
neighbours inside $G[N(u_{d-1})]$, which by Lemma~\ref{lem:matching} is a matching. Hence
$a(u_{d-1})\ge 2$, witnessed by $y = u_d$; the vertex $x$ is still a usable side-neighbour of
$u_0$ for the truncated geodesic, whose first step is still $u_1$; and Lemma~\ref{lem:g41} at
$d - 1 \ge 5$ gives $\pth(G) \ge (d-1)+4 = d+3$.
\end{proof}

The truncation was checked on twenty-four runs over graphs that do carry $a$-value $1$
vertices --- the leaf and triangle-leaf variants of the second witness named after
Theorem~\ref{thm:g42}, on $133$ and $134$ vertices, with $\ell = 764/133$ and $765/134$ and
radius $8$ --- since a truncation lemma checked only on graphs with $\mu \ge 2$ would be
checked exactly where it cannot be needed.

\begin{corollary}[one configuration remains]\label{cor:g44}
\REV{w133\_r11\_g43.py}
Let $G$ be connected and $C_4$-free with $\ell(G)>4$ and $\rad(G)\ge 5$. Then
$\pth(G)\ge \rad(G)+4$ unless every geodesic of length $\ge \rad(G)$ that fails to be
$3$-capped at an end has length exactly $\rad(G)$ \emph{and} $a$-value $1$ at its far end.
\end{corollary}

\begin{proof}
Some geodesic of length $\ge \rad(G)$ fails to be $3$-capped at an end, by
Corollary~\ref{cor:g38}. If one such has length $d \ge \rad(G)+1\ge 6$, then
Lemma~\ref{lem:g43} gives $\pth(G)\ge d+3\ge \rad(G)+4$. If one such has $a\ge 2$ at its far
end, then Lemma~\ref{lem:g41} gives $\pth(G)\ge d+4\ge \rad(G)+4$.
\end{proof}

\begin{remark}[the exceptional configuration is not exhibited]\label{rem:nonexhibited}
Corollary~\ref{cor:g44} reduces the slice $\rad \ge 5$ to a single configuration, and it is
important to say precisely how much is known about that configuration: it is \emph{not
exhibited}. What the campaign exhibits is the strictly broader class of $C_4$-free graphs with
$\ell > 4$, $\rad \ge 5$ and at least one vertex of $a$-value $1$, which is inhabited --- the
second witness of Theorem~\ref{thm:g42} with a leaf attached, and with a triangle-leaf
attached, are members. No graph is known to the campaign in which \emph{every} non-$3$-capped
geodesic of length $\ge \rad$ has length exactly $\rad$ with an $a$-value $1$ far end. So the
``unless'' clause of Corollary~\ref{cor:g44} could in principle be vacuous. That would be good
news; it is not knowledge, and until a witness exists the corollary must be read as reducing
the problem rather than as locating a live obstruction. The graph $Q$ of
Theorem~\ref{thm:g48} below is the nearest thing on file --- it has $\mu = 1$ and every one of
its $480$ cap-break frames has length exactly $\rad$ with far-end $a$-value $1$ --- but it is
not a witness for the exceptional configuration, because $\pth(Q)\ge 14 > \rad(Q)+4$, so the
``unless'' clause does no work there.
\end{remark}

\subsection{A refutation inside the residual}

One reformulation of the remaining problem, natural enough that the campaign pursued it as an
equivalent, is false.

\begin{theorem}[the far-end reformulation is false]\label{thm:g48}
\MV{w133\_r13\_q32\_verify.py}
It is not true that every connected $C_4$-free graph with $\ell > 4$ and $\rad \ge 5$ has a
geodesic of length $\ge \rad$ carrying a usable side-neighbour of $a$-value $\ge 4$ at one end
and $a$-value $\ge 2$ at the other.
\end{theorem}

\begin{proof}
Let $Q$ be the incidence graph of $\mathrm{PG}(2,4)$, on $42$ vertices, with a pendant path
$s - h_s - f_s$ of two new vertices attached at each of the six vertices given by
$[1{:}0{:}0]$, $[0{:}1{:}0]$ and $[0{:}0{:}1]$ read both as points and as lines. Recomputed
from that recipe rather than assumed: $Q$ has $n = 54$, $117$ edges, is connected,
triangle-free and $C_4$-free in the strong sense; its $a$-values are $6$ six times, $5$
thirty-six times, $2$ six times and $1$ six times, so $\sum_v a(v) = 234$ and
$\ell(Q) = 13/3 > 4$; and $\rad(Q) = 5$, $\diam(Q) = 7$. Over all triples $(u_0,u_1,u_d)$ with
$u_0$ and $u_d$ at distance $\ge \rad(Q)$, the number admitting a usable side-neighbour of
$a$-value $\ge 4$ at $u_0$ together with $a(u_d)\ge 2$ is zero. The count is not vacuous:
$480$ such triples exist with a usable side-neighbour of $a$-value $\ge 4$ at $u_0$, and every
one of them has length exactly $\rad(Q) = 5$ with far-end $a$-value exactly $1$.
\end{proof}

The same graph kills a second route. Peeling its six $a$-value $1$ vertices removes exactly
twelve vertices and returns the bare incidence graph of $\mathrm{PG}(2,4)$, on which $\ell$
stays above $4$ (it is $210/42 = 5$) while $\rad$ drops from $5$ to $3$. So no argument may
assume that peeling preserves the radius.

What survives is worth stating exactly, because a refutation is easy to over-read.
Theorem~\ref{thm:g48} does \emph{not} refute the underlying problem, which is to control the
radius drop under peeling or to trade it against the induced path the peeled material itself
supplies. On $Q$ that problem is comfortably satisfied: a depth-first search that makes no use
of the reformulation finds an induced path on at least $14$ vertices, against $\rad + 4 = 9$.
The live route is the quantitative exchange $\rad(G)\le \rad(G')+h$, where $G'$ is the peeled
graph and $h$ the length of the peeled hairs, and $Q$ is its tight control and an equality
instance: $h = 2$, $\rad(Q) = 5$, $\rad(Q') = 3$, and $5 \le 3+2$ holds with equality. Any
proof of that exchange has to execute on this graph. Finally, Theorem~\ref{thm:g42},
Lemma~\ref{lem:g43} and Corollary~\ref{cor:g44} are untouched by Theorem~\ref{thm:g48} and are
corroborated by it: they govern the $\mu \ge 2$ branch and the reduction to one configuration,
and $Q$ has $\mu = 1$ and is an instance of that configuration, so it shows the reduction is
not vacuous rather than threatening it.

\begin{table}[ht]
\centering
\small
\begin{tabular}{@{}lll@{}}
\toprule
Slice or route & Status & Where \\
\midrule
$\rad \ge 5$, $\mu \ge 2$ & closed & Thm.~\ref{thm:g42} \\
$\rad \ge 5$, $\mu = 1$ & \textbf{open}; one configuration & Lem.~\ref{lem:g43},
   Cor.~\ref{cor:g44} \\
\quad via the far-end reformulation & dead & Thm.~\ref{thm:g48} \\
\quad via ``peeling preserves the radius'' & dead & Thm.~\ref{thm:g48} \\
\quad via the exchange $\rad(G)\le \rad(G')+h$ & live, tight-controlled by $Q$ &
   Thm.~\ref{thm:g48} \\
$\rad \le 4$ & \textbf{open} & \S\ref{sec:open} \\
\bottomrule
\end{tabular}
\caption{The band $\lfloor\ell\rfloor\ge 4$ after this section. Nothing here closes the band.}
\label{tab:pocket2}
\end{table}

So the residual of the whole band is the single configuration of Corollary~\ref{cor:g44} at
$\rad \ge 5$ --- every non-$3$-capped geodesic of length $\ge \rad$ having length exactly
$\rad$ and an $a$-value $1$ far end, a configuration no known graph realises --- together with
the radii $2$, $3$ and $4$, of which radius $4$ additionally wants the $a(u_0)\ge 3$ upgrade
that Lemma~\ref{lem:g36} needs there.

\begin{remark}[computational provenance of this section]\label{rem:pocket2prov}
Every invariant of both witnesses has been computed twice, by implementations that share no
code. For $W_2$ the two implementations use different models of $\mathrm{PSL}(2,11)$, as
described in Proposition~\ref{prop:w2}. For $Q$ the first implementation,
\texttt{w133\_r13\_q32\_verify.py}, builds $\mathrm{PG}(2,4)$ from $\mathrm{GF}(4)$
arithmetic, normalising projective triples and testing incidence by the vanishing of the
bilinear form; the second, \texttt{w133\_r18\_q32\_rebuild.py}, builds it from a Singer
planar difference set, taking the points to be the residues modulo $21$, the lines to be the
$21$ translates of $D = \{0,1,6,8,18\}$, and incidence to be $p - i \in D$. No field, no
coordinates and no line of code are shared between the two, and the difference-set property is
verified rather than assumed. Both rebuild the graph from the recipe and assert
$C_4$-freeness before reading any statistic off it, and all of $n$, the degree sequence,
triangle-freeness, $\sum_v a(v)$, the $a$-profile, $\rad$, $\diam$, the peeling behaviour and
the induced-path bound agree. Cayley-type and incidence-type constructions are determined up
to isomorphism by their defining data and every quantity listed is an isomorphism invariant,
so the agreement is genuine independent confirmation rather than a re-run. The review tier
carried by Lemma~\ref{lem:g41} to Corollary~\ref{cor:g44} is nevertheless a statement about
human-style reading of proofs by adversarial reviewers, and it does not extend to these
computations at all.
\end{remark}

\section{Negative results, refutations, and named obstacles}\label{sec:limits}

This section is part of the report, not an appendix to it: each item closes a route, and
three of them close routes that the campaign itself had been pursuing.

\subsection{No two-endpoint lemma reaches $\rad + k$ for $k \ge 4$}

\begin{proposition}\label{prop:tm}
\MV{wowii133\_gplus.c}
For $m \ge 1$ let $T_m$ be the tree consisting of a path $u - s - v$ with $m$ pendant vertices
attached to $u$ and $m$ to $v$. Then $T_m$ is $C_4$-free with $n = 2m+3$, $\rad(T_m) = 2$,
$a(u) = a(v) = m+1$, $\ell(T_m) = (4m+4)/(2m+3) \to 2$, and $\pth(T_m) = 5$ for every $m$.
Consequently, for every $k \ge 4$ the statement
\begin{quote}
``$C_4$-free with a pair at distance $\rad$ whose $a$-values are both $\ge k$ implies
$\pth \ge \rad + k$''
\end{quote}
is false, and no strengthening of Lemma~\ref{lem:g1} that reads only the two endpoints of a
geodesic can prove (C133-b) in the band $\lfloor \ell\rfloor \ge 4$.
\end{proposition}

\begin{proof}
$T_m$ is a tree, hence $C_4$-free. $\ecc(s) = 2$ and $\ecc(u) = 3$, so $\rad = 2$ and
$\mathrm{dist}(u,v) = 2 = \rad$. The $a$-values are $m+1$ at $u$ and at $v$, $2$ at $s$ and
$1$ at each of the $2m$ pendant vertices, giving $\sum a = 4m+4$. A longest induced path is
$\text{leaf} - u - s - v - \text{leaf}$, on $5$ vertices. For $k \ge 4$ and $m \ge k-1$ the
hypothesis holds while $\rad + k = 2+k \ge 6 > 5 = \pth$.
\end{proof}

Note that $T_m$ does \emph{not} contradict the conjecture: $\lfloor \ell\rfloor = 2$ and
$\pth = 5 \ge \rad + 2$. What it kills is a proof technique. It also shows why
Theorem~\ref{thm:gplus} is stated with the hypothesis $\mu \ge 2$ rather than with a
hypothesis on the endpoints only: $\mu(T_m) = 1$.

\subsection{The layered-counting route dies above radius 2}

Theorem~\ref{thm:g7} proves the statement $(GP_2)$: \emph{$C_4$-free with $\rad = 2$ and
$\ell \ge 3$ implies there is a pair at distance $\ge \rad$ with $a$-values
$(\ge 3, \ge 2)$}. The obvious next step is $(GP_3)$.

\begin{proposition}\label{prop:gp3}
\MV{w133\_r3\_counterexample.py}
$(GP_3)$ is false. Let $A$ be the Hoffman--Singleton graph with two pendant vertices attached
to every vertex. Then $A$ is connected and $C_4$-free with $n = 150$, $m = 275$, $\rad = 3$,
$\diam = 4$, $\sum_v a(v) = 550$, $\ell = 11/3$ and $\lfloor \ell\rfloor = 3$ --- so $A$ lies
in the band under attack --- while the number of pairs at distance $\ge \rad$ with $a$-values
$(\ge 3,\ge 2)$ is \emph{zero}: the largest distance between two vertices of $a$-value
$\ge 2$ is $2 < 3 = \rad$. The graph $B$ obtained by attaching only three pendant vertices,
at three mutually distance-$2$ vertices with no common neighbour, has $n = 53$, $\rad = 3$,
$\lfloor\ell\rfloor = 6$ and again zero such pairs. Both satisfy the conjecture ($A$ has an
induced path on $\ge 16$ vertices against a requirement of $6$).
\end{proposition}

\begin{proof}[Verification]
The vertices of $a$-value $\ge 2$ are exactly the $50$ Hoffman--Singleton vertices, which lie
pairwise within distance $\diam(\mathrm{HS}) = 2$; every vertex at distance $\ge 3$ from
anything is a pendant vertex, of $a$-value $1$. The invariants are computed by the named
script.
\end{proof}

The consequence is sharper than it looks. The route that proves Theorem~\ref{thm:g7} aims to
show ``$(\star)$ and $\rad \ge 3$ and a high vertex exists $\Rightarrow 3n - \sum a > 0$'';
the witness $A$ satisfies $(\star)$, has $\rad = 3$ and $50$ high vertices, and has
$3n - \sum a = 450-550 = -100 < 0$. So the fibre count of Theorem~\ref{thm:g7} cannot be
generalised to radius $3$, and any proof there must genuinely construct an induced path
rather than count. Only three pendant vertices suffice to lift the radius from $2$ to $3$ and
destroy every good pair.

\subsection{The deep endpoint lemma cannot be pushed to $d = 3$}

Lemma~\ref{lem:g36} extends Lemma~\ref{lem:g14} to $d = 4$ with no extra hypothesis, and
Lemma~\ref{lem:g37} extends it to $d = 3$ under the hypothesis that some usable side-neighbour
$y$ of $u_3$ satisfies $y \not\sim x$. The complementary configuration --- call a frame
\emph{forced} if every usable side-neighbour $y$ of $u_3$ satisfies $y \sim x$ --- is not
merely unproved. It is false, twice over.

\begin{proposition}[the bare statement is false]\label{prop:ce1}
\MV{w133\_r10\_adjudicate.py}
Let $G_1$ be the graph on $\{0,\dots,9\}$ with edges
\[
  09,\; 05,\; 15,\; 23,\; 25,\; 26,\; 28,\; 49,\; 46,\; 57,\; 58 .
\]
Then $G_1$ is connected and $C_4$-free; $2-6-4-9$ is a geodesic of length $3$;
$a(2) = 3$ and $a(9) = 2$; $x = 5$ is a usable side-neighbour of $u_0 = 2$ with $a(5) = 4$;
the frame is forced (the unique usable side-neighbour of $u_3 = 9$ is $0$, and $0 \sim 5$);
and $\pth(G_1) = 6 < 7 = d+4$. Hence the $d = 3$ analogue of Lemma~\ref{lem:g36} is false as
stated.
\end{proposition}

The vertices of $a$-value $1$ in $G_1$ are the leaves $1,3,7$ and the triangle-leaf $8$;
peeling them leaves a $6$-cycle, in which the frame no longer exists. This is what made the
following ``$F4$ form'' --- the same statement with the extra hypothesis $\mu(G)\ge 2$, which
is free in the intended application because peeling preserves $C_4$-freeness and raises the
slack (Lemma~\ref{lem:g11}) --- look like the right repair. It is not.

\begin{theorem}[the repaired statement is also false]\label{thm:g45}
\MV{w133\_r11\_adjudicate.py}
Let $H$ be the graph on $\{0,\dots,9\}$ with edges
\[
  01,\;04,\;05,\;08,\;12,\;23,\;26,\;35,\;37,\;39,\;46,\;47,\;48,\;89 .
\]
Then $H$ is connected and $C_4$-free with $a$-value vector
$(3,2,3,4,3,2,2,2,2,2)$, so $\mu(H) = 2$ and \emph{no peeling is possible};
$2-1-0-8$ is a geodesic of length $3$ with $a(u_0) = 3$ and $a(u_3) = 2$; $x = 3$ is a usable
side-neighbour of $u_0$ with $a(x) = 4$; the frame is forced, the unique usable
side-neighbour of $u_3 = 8$ being $9$, with $9 \sim 3$; and $\pth(H) = 6 < 7$, witnessed by
the induced path $0-8-9-3-2-6$ and by exhaustive search for the upper bound.
\end{theorem}

\begin{proof}[Verification, from the edge list alone]
$N(0) = \{1,4,5,8\}$ with the single internal edge $48$, so $a(0) = 3$; $N(2) = \{1,3,6\}$
is independent, so $a(2) = 3$ and every neighbour of $2$ outside $u_1 = 1$'s component is
usable; $N(3) = \{2,5,7,9\}$ is independent, so $a(3) = 4$; $N(8) = \{0,4,9\}$ with the
single internal edge $04$, so $a(8) = 2$ and the usable side-neighbour set of $8$ with
respect to the geodesic is exactly $\{9\}$, and $39$ is an edge. No two vertices have two
common neighbours. Neither $2$ nor $8$ has a neighbour in common with the other, and
$2-1-0-8$ is a shortest path. An exhaustive enumeration of induced paths returns $6$ as the
maximum.
\end{proof}

\begin{remark}[scope, and why $G_1$ and $H$ do not touch \S\ref{sec:adopted}]
$\sum_v a(v) = 25$ and $n = 10$ give $\ell(H) = 2.5 < 3$, so $H$ lies \emph{outside} the live
class of \S\ref{sec:peeling}, which carries $\ell > 3$. What $H$ refutes is therefore the
hypothesis-free \emph{local} route at $d = 3$ --- the deep endpoint extension cannot be
pushed down to $d = 3$ in this currency, with or without $\mu \ge 2$ --- and not
Problem~\ref{prob:rpd}. The next attempt at $D = 3$ must consume either a global hypothesis
such as $\ell > 3$ or a declared $a$-value surcharge. Neither Lemma~\ref{lem:g36} nor
Lemma~\ref{lem:g37} is touched: $\diam(H) = 3$, so $H$ carries no geodesic of length $\ge 4$
and Lemma~\ref{lem:g36} cannot execute; and an exhaustive enumeration of $H$'s qualifying
$d = 3$ frames returns exactly two, $2-1-0-8$ and $2-6-4-8$, both forced, so
Lemma~\ref{lem:g37} cannot execute either.
\end{remark}

\begin{remark}[a methodological correction]\label{rem:frames}
Theorem~\ref{thm:g45} was found by the adversarial reviewer assigned to \S\ref{sec:adopted},
on a target that the campaign's own sweep had recorded as open after finding no
counterexample among $1112$ \emph{forced frames} (of which $63$ satisfied $\mu \ge 2$). The
sweep was frame-local: it enumerated local skeletons and could not see $\pth(G)$ of a
completed graph. The lesson, recorded here because the same species of error appears twice in
the campaign's logs, is that ``$0$ counterexamples in $N$ frames'' is evidence about frames
and never about graphs, and must not be written in a way that reads otherwise.
\end{remark}

\subsection{Named obstacles}

\begin{proposition}[surgery]\label{prop:o1}
\MV{wowii133\_families.py}
Let $G^\ast$ be the Hoffman--Singleton graph with one edge subdivided. Then $G^\ast$ is
$C_4$-free with $n = 51$, $\delta = 2$, $\rad = 3$, $\ell = 6.902$ and
$\lfloor\ell\rfloor = 6$, so the conjecture demands $\pth \ge 9$ (true: $\pth \ge 20$) while
$\max(\rad+2,\ \delta+2,\ \rad+\lfloor\delta/2\rfloor) = 5$. Hence Theorems~\ref{thm:D},
\ref{thm:E} and \ref{thm:F} cannot be combined into a proof of the conjecture.
\end{proposition}

Every bound in \S\ref{sec:closed} is driven by $\delta$ or $\rad$; both are destroyed by one
local modification while $\ell$ moves by $O(1/n)$. A correct proof must use a quantity as
robust as $\ell$.

Three further obstacles are recorded without proof, since each is an observation rather than
a theorem. \emph{(O2) The degeneracy factor $2$.} Every greedy argument really bounds $\pth$
by $2 + \max_{H\subseteq G}\delta(H)$, i.e.\ by the degeneracy, which can be as small as half
the average degree, whereas the conjecture needs about the average degree; closing that
factor for $C_4$-free graphs is a sharp form of the known results on induced subdivisions in
$K_{s,s}$-free graphs of large average degree, for which only tower-type bounds are
available. (This attribution is from memory: the campaign ran without web access and it was
not checked against the literature.) \emph{(O3) No per-vertex charging.} The star $K_{1,m}$
has $a(\text{centre}) = m$ but $\pth = 3$, so no bound of the form
``$\pth \ge f(a(v)) + g(v)$ for some vertex $v$'' can exist: any proof must genuinely
average. \emph{(O4) Sharpness at the bottom.} By Proposition~\ref{prop:tight} five graphs
attain equality, so no argument may lose even one vertex.

Finally, one negative result that saves computation:

\begin{proposition}\label{prop:circulant}
\PRVn
Every circulant graph of degree $\ge 4$ contains a $4$-cycle, since a connection set
containing $\pm a$ and $\pm b$ with $a \ne b$ gives $0 = a + (-a) = b + (-b)$ and hence two
vertices with two common neighbours. Circulants are therefore useless as test samples for any
$C_4$-free statement with a degree floor $\ge 4$.
\end{proposition}

\section{Status and open problems}\label{sec:open}

The conjecture is open. What remains, in the order in which the campaign ranks it:

\begin{enumerate}
\item \textbf{Problem~\ref{prob:rpd} at $D = 3$.} By Theorem~\ref{thm:r3}, a minimum-order
counterexample in the band $\lfloor\ell\rfloor = 3$ has radius $\ge 4$, and
\S\ref{sec:peeling} reduces radius $\ge 4$ to Problem~\ref{prob:rpd} at $D = \rad - 1 \ge 3$.
The case $D = 2$ is Theorem~\ref{thm:g12}; $D = 4$ and $D \ge 5$ are available from
Lemma~\ref{lem:g36} and Lemma~\ref{lem:g14} \emph{under the side-neighbour hypothesis}
$a(x)\ge 4$, which still has to be discharged from $\ell > 3$. The case $D = 3$ is the first
genuinely open one, and Theorem~\ref{thm:g45} says the local route to it is closed: a global
hypothesis must be paid for.

\item \textbf{Discharging the side-neighbour hypothesis.} Corollary~\ref{cor:g38} discharges
exactly this hypothesis, but from $\ell > 4$ rather than $\ell > 3$. An analogue at
$\ell > 3$ would combine with Lemmas~\ref{lem:g14} and \ref{lem:g36} to give
Problem~\ref{prob:rpd} for $D = 4$ and $D \ge 5$ outright.

\item \textbf{The band $\lfloor\ell\rfloor \ge 4$.} Corollary~\ref{cor:g38} supplies a
peripheral geodesic with a high side-neighbour at every radius; what a proof still needs at
that geodesic is $a(u_0)\ge 3$ and $a(u_d)\ge 2$. At $\rad \ge 5$ the first of these is
free (Lemma~\ref{lem:g41}) and \S\ref{sec:pocket2} reduces the band to Table~\ref{tab:pocket2}:
the surviving route is the exchange $\rad(G)\le \rad(G')+h$ under peeling, whose tight
control is the graph $Q$ of Theorem~\ref{thm:g48}, together with the radii $\le 4$, radius
$4$ additionally wanting the $a(u_0)\ge 3$ upgrade. The most informative single thing that
could be done next is to exhibit --- or rule out --- a graph realising the exceptional
configuration of Corollary~\ref{cor:g44}; see Remark~\ref{rem:nonexhibited}.

\item \textbf{A7.} Proposition~\ref{prop:aux} records the real-valued strengthening
$\ell(G)\le \pth(G)-\rad(G)$, verified over the whole exhaustive range and tight on the same
five graphs. Proving it would prove the conjecture, and Theorem~\ref{thm:g12} is a proved
non-trivial instance of it. It remains the campaign's recommended target.

\item \textbf{Formalisation.} Nothing here is formalised. The one convention that a Lean
formalisation must not lose is the strong reading of $C_4$-freeness (\S\ref{sec:intro}):
several lemmas above --- Lemma~\ref{lem:g23}(a) among them --- are false under the
induced-only reading.
\end{enumerate}

\appendix

\section{Verification artifacts}\label{app:artifacts}

Every script named below lives in the campaign repository at the stated relative path; each
asserts the statements attributed to it and exits non-zero on the first violation. Logs are
the same filename with extension \texttt{.out}. No result in this paper was obtained with a
SAT solver.

\begin{longtable}{@{}p{0.32\textwidth}p{0.62\textwidth}@{}}
\toprule
\textbf{Artifact} & \textbf{What it certifies} \\
\midrule
\endfirsthead
\toprule
\textbf{Artifact} & \textbf{What it certifies} \\
\midrule
\endhead
\texttt{notes/proofs/} \texttt{wowii133\_verify.c} &
The exhaustive generator and assertion harness: all connected graphs on $n \le 10$
($136\,362\,861$ evaluations) and all connected $C_4$-free graphs on $n\le 13$
($13\,544\,784$); asserts Theorems~\ref{thm:C}, \ref{thm:D}, \ref{thm:E}, \ref{thm:F} and the
coverage census of \S\ref{sec:falsification}; self-validates its class counts against the
number of connected graphs on $n$ nodes. \\
\texttt{notes/proofs/} \texttt{wowii133\_families.py} &
The $79$ structured graphs up to $n = 183$ (Petersen, Hoffman--Singleton, incidence graphs of
$PG(2,q)$, polarity graphs $ER_q$, cages, random girth-$5$ regular graphs), their invariants,
and the surgery witness of Proposition~\ref{prop:o1}. \\
\texttt{notes/proofs/} \texttt{wowii133\_sa.c} &
The adversarial simulated-annealing search over connected $C_4$-free graphs,
$n = 10..24$, minimising $\pth-\rad-\lfloor\ell\rfloor$ ($\approx 2\cdot10^6$ proposals,
minimum slack $1$ for $n \ge 14$). \\
\texttt{problems/wowii/} \texttt{wowii133\_gplus.c} &
The same harness extended with assertions for Lemmas~\ref{lem:g10}, \ref{lem:g1},
\ref{lem:g9} (including the $334$ failures of the $d=3$ variant), \ref{lem:g2},
Theorem~\ref{thm:gplus}, Lemma~\ref{lem:g13}, the structural steps of
Theorem~\ref{thm:g12}, and the hexagon-rigidity assertions ($1272$ non-vacuous induced-$C_6$
instances); the coverage census $606 \to 0$. \\
\texttt{problems/wowii/} \texttt{w133\_r3\_counterexample.py} &
Proposition~\ref{prop:gp3} (witnesses $A$ and $B$) and the residual analysis of
\S\ref{sec:peeling}; the frames backing Lemma~\ref{lem:g14}. \\
\texttt{problems/wowii/} \texttt{w133\_g12\_nonempty.py} &
Proposition~\ref{prop:g12nonempty}, triangle-free branch: Hoffman--Singleton lies in
Theorem~\ref{thm:g12}'s hypothesis class, with an induced-path certificate. \\
\texttt{problems/wowii/} \texttt{w133\_res3\_search.py} &
Proposition~\ref{prop:g12nonempty}, triangle branch ($ER_5,ER_7,ER_{11}$); the nine witnesses
of Remark~\ref{rem:live} and the verification that all nine violate (R3b). \\
\texttt{problems/wowii/} \texttt{w133\_res\_kill.py} &
Remark~\ref{rem:r3a} (the five members of $\mathrm{RES(a)}\setminus\mathrm{RES}$);
Lemmas~\ref{lem:mass}, \ref{lem:g18}, \ref{lem:g19}, \ref{lem:g4}, \ref{lem:g5},
Theorem~\ref{thm:g7} over a $48$-graph pool ($1363$ distance-$3$ pairs, $329$ subsets, $371$
eccentricity-$2$ vertices, $301$ centres). \\
\texttt{problems/wowii/} \texttt{w133\_r6\_adjudicate.py} &
Lemma~\ref{lem:g22} for $k = 3,4,5$ ($4226 / 67\,344 / 194\,280$ endpoints) and the
friendship-graph tightness; the star counterexample that establishes the gap repaired by
Step~1 of Theorem~\ref{thm:g12}; the classification of triangle-free diameter-$2$ $C_4$-free
graphs on $n \le 7$ as $25$ stars plus $12$ Moore graphs. \\
\texttt{problems/wowii/} \texttt{w133\_r7\_triend.py} &
Lemma~\ref{lem:g25} and the layer decomposition of Lemma~\ref{lem:g23} on $79\,852$
triangular-endpoint carriers, plus the two bounded integer sweeps used in
Theorem~\ref{thm:g26}. \\
\texttt{problems/wowii/} \texttt{w133\_r8\_q61.py} &
Lemmas~\ref{lem:g32p}, \ref{lem:g32}, Corollary~\ref{cor:g32.1} ($1330$ hexagons, $1300$
classified off-hexagon vertices, $15\,960$ mass identities); Lemma~\ref{lem:g33}'s $176$
frames; the exhaustive frame enumeration behind Theorem~\ref{thm:g35}, whose realised maximum
mass is exactly $6$. \\
\texttt{problems/wowii/} \texttt{w133\_r9\_frame\_l.py} &
The five admissible frames of Remark~\ref{rem:secondkill} with $\ell$ computed;
$\max \ell = 8/3$. \\
\texttt{problems/wowii/} \texttt{w133\_r9\_controls.py} &
The verified control graphs: Petersen, the $PG(3,2)$ incidence graph ($n=50$, $\ell = 4.2$,
$\rad = 3$) and the $PG(2,3)$ incidence graph ($n = 26$, $\ell = 4.0$), with exact
induced-path search. \\
\texttt{problems/wowii/} \texttt{w133\_r10\_adjudicate.py} &
Proposition~\ref{prop:ce1}; the $d=3$ and $d=4$ frame sweeps; the lemma chain of
Theorem~\ref{thm:g38} on $449$ peripherally $3$-capped graphs; the verification that
$PG(2,3)$ and $PG(3,2)$ are not peripherally $3$-capped. \\
\texttt{problems/wowii/} \texttt{w133\_r11\_adjudicate.py} &
Theorem~\ref{thm:g45} re-verified from the edge list by independent code, including
$\pth(H) = 6$ by exhaustion, $\ell(H) = 2.5$, and the enumeration of $H$'s two qualifying
$d = 3$ frames; Remark~\ref{rem:b10}'s peeling arithmetic on $G_1$. \\
\texttt{problems/wowii/} \texttt{w133\_r11\_qradp.py} &
The control $W_2$ of Proposition~\ref{prop:w2} in the matrix model of
$\mathrm{PSL}(2,11)$, and the corresponding larger members of the family. \\
\texttt{problems/wowii/} \texttt{w133\_r17\_w2\_recompute.py} &
The independent recomputation of $W_2$: $\mathrm{PSL}(2,11)$ realised as a permutation group
on the projective line over $\mathrm{GF}(11)$, sharing no code with the above; faithfulness
checked, eccentricity $5$ verified at all $660$ vertices, and a certified induced path on
$192$ vertices. Its companion \texttt{w133\_r17\_v7probe.py} settles which of its own
assertions are refutable by injecting a fault against each. \\
\texttt{problems/wowii/} \texttt{w133\_r11\_g14sharp.py} &
Lemma~\ref{lem:g41} on $120$ frames with $a(u_0) = 2$ exactly --- that is, non-vacuously
where the deleted hypothesis fails. \\
\texttt{problems/wowii/} \texttt{w133\_r11\_g42.py} &
Theorem~\ref{thm:g42}'s mechanism run end to end $50$ times on two witnesses, including the
constructed path of $9 = \rad+4$ vertices on $W_2$ of Remark~\ref{rem:twonumbers}. \\
\texttt{problems/wowii/} \texttt{w133\_r11\_g43.py} &
Lemma~\ref{lem:g43} on $24$ truncation runs over graphs that carry $a$-value $1$ vertices. \\
\texttt{problems/wowii/} \texttt{w133\_r13\_q32\_verify.py} &
The graph $Q$ of Theorem~\ref{thm:g48} rebuilt from its recipe, $22$ checks: the invariants,
the empty count of qualifying triples against the non-empty count of $480$ cap-break frames,
and the radius drop $5 \to 3$ under peeling. \\
\texttt{problems/wowii/} \texttt{w133\_r18\_q32\_rebuild.py} &
The same graph $Q$ rebuilt independently, $37$ checks, with $\mathrm{PG}(2,4)$ obtained from
a Singer planar difference set rather than from $\mathrm{GF}(4)$ arithmetic; it reproduces
every invariant, including the count of $480$ cap-break frames exactly, and injects a fault
against each of its four load-bearing certifiers to show they can fail. \\
\texttt{problems/wowii/} \texttt{w133\_r15\_key.py} &
The held-out key for the second adversarial pass over \S\ref{sec:pocket2}: the witnesses'
invariants recomputed from scratch, $42$ checks. \\
\texttt{prompts/} \texttt{w133\_S3\_R24.md}, \texttt{w133\_S3\_R24\_v6\_qwen.md} &
The two self-contained briefs placed under the two adversarial passes that
Lemmas~\ref{lem:g41}, \ref{lem:g43}, Theorem~\ref{thm:g42} and Corollary~\ref{cor:g44}
carry: the statements and full proofs exactly as reviewed, each with its held-out table and
its named joints. The reviewers' reports are retained in the campaign's quarantine directory
for untrusted-channel output and were never copied into the repository. \\
\texttt{problems/wowii/} \texttt{w133\_r8\_judge\_*.py} &
The first reviewer's own verification code for \S\ref{sec:tri}, written independently of the
campaign's: an exhaustive pool of all $4448$ connected $C_4$-free graphs on $n\le 10$ whose
class counts independently reproduce the campaign's, plus explicit families;
$\approx 90\,000$ assertions, $0$ failures. \\
\texttt{problems/wowii/} \texttt{w133\_r8\_judge20\_*.py} &
The blind reviewer's own verification code for \S\ref{sec:hexagon}, including the
counterexample of Remark~\ref{rem:g32false} and a re-verification of the repaired
Lemma~\ref{lem:g32} over $2\,105\,408$ frames. \\
\texttt{problems/wowii/} \texttt{w133\_S3\_R23\_opus.md} &
The adversarial review report on \S\ref{sec:adopted}, containing the refutation
$H$ of Theorem~\ref{thm:g45} and the peeling-arithmetic defect of Remark~\ref{rem:b10}. \\
\texttt{prompts/} \texttt{w133\_S3\_R23.md} &
The self-contained brief placed under that review: the statements and full proofs of
Lemmas~\ref{lem:g36}, \ref{lem:g37} and Theorem~\ref{thm:g38} exactly as reviewed, with the
named joints and the mandatory probes. \\
\texttt{problems/wowii/} \texttt{w133\_S3\_GDIFF\_qwen.md} &
The diff-scoped second review pass over \S\ref{sec:hexagon} and the repair of
Remark~\ref{rem:g32false}, returning no mathematics defect. \\
\texttt{problems/wowii/} \texttt{w133\_S3\_G12\_qwen.md} &
The review pass that located the gap in Theorem~\ref{thm:g12} Step~1, repaired there. \\
\texttt{notes/proofs/} \texttt{wowii133\_draft.md} &
The campaign's running attack report: the full round-by-round record, including every
adjudication, every retraction, and the material omitted from this paper. \\
\bottomrule
\end{longtable}

\end{document}
